\documentclass[12pt]{article}

\usepackage[T1]{fontenc}
\usepackage[utf8]{inputenc}
\usepackage[brazil]{babel}
\usepackage{lmodern}
\usepackage{microtype}
\usepackage[a4paper,margin=2.5cm]{geometry}
\usepackage{setspace}
\usepackage{indentfirst}
\usepackage{amsmath}

\setstretch{1.15}

\title{Uma Estrutura Teórica para a Análise Monetária \thanks{Texto Traduzido livremente por Luiz Mario Andrade com o auxílio da ferramenta Chat GPT 5.4}}
\author{Milton Friedman}
\date{Março – Abril, 1970}

\begin{document}

\maketitle

Todo estudo empírico repousa sobre uma estrutura teórica, sobre um conjunto de hipóteses tentativas que a evidência é desenhada para testar ou adumbrar. Pode ajudar o leitor da série de monografias sobre a moeda que Anna J. Schwartz e eu temos escrito para expor explicitamente a estrutura teórica geral que as subjaz.\footnote{Vários revisores de nossa obra \textit{A Monetary History of the United States, 1867–1960} (Friedman e Schwartz 1963b) nos criticaram por não tornar explícita a estrutura teórica empregada naquele livro. Este artigo é, em grande parte, uma resposta a essa crítica. Ver Culbertson (1964) e Meltzer (1965).}

Essa estrutura é a teoria quantitativa da moeda — uma teoria que assume muitas formas diferentes e remonta ao próprio início do pensamento sistemático sobre questões econômicas. Ela provavelmente foi "testada" com dados quantitativos de forma mais extensiva do que qualquer outro conjunto de proposições na economia formal — a menos que seja a curva de demanda de inclinação negativa. No entanto, a teoria quantitativa tem sido um pomo de discórdia contínuo. Até as últimas três décadas, ela era geralmente apoiada por estudantes sérios de economia, aqueles que hoje chamaríamos de economistas profissionais, e rejeitada por leigos. No entanto, o sucesso da revolução keynesiana levou à sua rejeição por talvez a maioria dos economistas profissionais. Só recentemente ela experimentou um renascimento, de modo que mais uma vez comanda a adesão de muitos economistas profissionais. Tanto sua aceitação quanto sua rejeição basearam-se basicamente em julgamentos sobre regularidades empíricas.

\section{A Teoria Quantitativa: Quantidade Nominal versus Real de Moeda}

Em todas as suas versões, a teoria quantitativa repousa sobre uma distinção entre a quantidade \textit{nominal} de moeda e a quantidade \textit{real} de moeda. A quantidade nominal de moeda é a quantidade expressa em quaisquer unidades usadas para designar a moeda — talentos, siclos, libras, francos, liras, dracmas, dólares e assim por diante. A quantidade real de moeda é a quantidade expressa em termos do volume de bens e serviços que a moeda pode comprar.

Não há uma maneira única de expressar a quantidade real de moeda. Uma maneira de expressá-la é em termos de uma cesta padrão especificada de bens e serviços. É isso que implicitamente se faz quando a quantidade real de moeda é calculada dividindo-se a quantidade nominal de moeda por um índice de preços. A cesta padrão é então a cesta cujos componentes são usados como pesos no cálculo do índice de preços — geralmente, a cesta comprada por algum grupo representativo em um ano-base.

Uma maneira diferente de expressar a quantidade real de moeda é em termos das durações de tempo dos fluxos de bens e serviços que a moeda poderia comprar. Para uma família, por exemplo, a quantidade de moeda pode ser expressa em termos do número de semanas do nível médio de consumo da família que ela poderia financiar com seus saldos monetários, ou, alternativamente, em termos do número de semanas de sua renda média a que seus saldos monetários são iguais. Para uma empresa comercial, a quantidade real de moeda que ela detém pode ser expressa em termos do número de semanas de suas compras médias, ou de suas vendas médias, ou de suas despesas médias com serviços produtivos finais (valor líquido adicionado) a que seus saldos monetários são iguais. Para a comunidade como um todo, a quantidade real de moeda pode ser expressada em termos do número de semanas das transações agregadas da comunidade, ou do produto líquido agregado da comunidade, ao qual ela é igual.

O recíproco de qualquer uma desta última classe de medidas da quantidade real de moeda é uma velocidade de circulação para a unidade correspondente ou grupo de unidades. Em cada caso, o cálculo da quantidade real de moeda ou da velocidade é feito ao conjunto de preços vigentes na data a que o cálculo se refere. Esses preços são a ponte entre a quantidade nominal e a real de moeda.

A teoria quantitativa da moeda toma como certo que o que realmente importa para os detentores de moeda é a quantidade real, em vez da quantidade nominal que eles detêm, e que existe uma quantidade real de moeda razoavelmente definida que as pessoas desejam manter sob quaisquer circunstâncias dadas. Suponha que a quantidade nominal que as pessoas detêm em um momento particular por acaso corresponda, aos preços correntes, a uma quantidade real maior do que a quantidade que desejam manter. Os indivíduos buscarão então se desfazer do que consideram saldos monetários excessivos; eles tentarão pagar uma soma maior para a compra de títulos, bens e serviços, para o pagamento de dívidas e como presentes do que estão recebendo das fontes correspondentes. No entanto, eles não podem, como grupo, ter sucesso. As despesas de um homem são as receitas de outro. Um homem pode reduzir seu saldo monetário nominal apenas persuadindo outra pessoa a aumentar o seu. A comunidade como um todo não pode, em geral, gastar mais do que recebe.

A tentativa de fazê-lo terá, no entanto, efeitos importantes. Se os preços e a renda estiverem livres para mudar, a tentativa de gastar mais elevará o volume de despesas e receitas, expresso em unidades nominais, o que levará a uma licitação de preços e talvez também a um aumento no produto. Se os preços forem fixados por costume ou por decreto governamental, a tentativa de gastar mais será acompanhada por um aumento em bens e serviços ou produzirá "escassez" e "filas". Estas, por sua vez, elevarão o preço efetivo e provavelmente, mais cedo ou mais tarde, forçarão mudanças nos preços oficiais.

O excesso inicial de saldos nominais tenderá, portanto, a ser eliminado, embora não haja mudança na quantidade nominal de moeda, seja por uma redução na quantidade real disponível para manter através de aumentos de preços ou por um aumento na quantidade real desejada através de aumentos de produto. E inversamente para uma deficiência inicial de saldos nominais.

Fica claro desta discussão que mudanças nos preços e na renda nominal podem ser produzidas tanto por mudanças nos saldos reais que as pessoas desejam manter quanto por mudanças nos saldos nominais disponíveis para que as mantenham. De fato, é uma tautologia, resumida na famosa equação quantitativa, que todas as mudanças na renda nominal podem ser atribuídas a uma ou outra — assim como uma mudança no preço de qualquer bem pode sempre ser atribuída a uma mudança na demanda ou na oferta. No nível analítico, contudo, a teoria quantitativa é uma análise dos fatores que determinam a quantidade de moeda que a comunidade deseja manter; no nível empírico, é a generalização de que as mudanças nos saldos reais desejados (na demanda por moeda) tendem a ocorrer lenta e gradualmente ou a ser o resultado de eventos postos em marcha por mudanças prévias na oferta, enquanto que, em contraste, mudanças substanciais na oferta de saldos nominais podem e frequentemente ocorrem independentemente de quaisquer mudanças na demanda. A conclusão é que mudanças substanciais nos preços ou na renda nominal são quase invariavelmente o resultado de mudanças na oferta nominal de moeda.

\section{Equações Quantitativas}

A tautologia incorporada na equação quantitativa é um dispositivo útil para esclarecer as variáveis enfatizadas na teoria quantitativa. A equação quantitativa assumiu diferentes formas, de acordo com as diferentes variáveis que os teóricos quantitativos enfatizaram.

\subsection{Equação de Transações}

A versão mais famosa da equação quantitativa é, sem dúvida, a versão de transações popularizada por Irving Fisher (Fisher 1911, pp. 24–54):
\begin{equation}
MV = PT
\end{equation}
ou
\begin{equation}
MV + M'V' = PT.
\end{equation}

Nesta versão, o evento elementar é uma transação: uma troca na qual um ator econômico transfere a outro bens ou serviços ou títulos e recebe em troca uma transferência de moeda. O lado direito das equações corresponde à transferência de bens, serviços e títulos; o lado esquerdo, à transferência correspondente de moeda.

Cada transferência de bens, serviços ou títulos é considerada o produto de um preço e uma quantidade: salário por semana vezes o número de semanas, preço de um bem vezes o número de unidades do bem, dividendo por ação vezes o número de ações, preço por ação vezes o número de ações, e assim por diante. O lado direito das equações (1) e (2) é o agregado de tais pagamentos durante algum intervalo, com $P$ uma \textit{média} adequadamente escolhida dos preços, e $T$ um \textit{agregado} adequadamente escolhido das quantidades durante esse intervalo, de modo que $PT$ seja o valor nominal total dos pagamentos durante o intervalo em questão. As unidades de $P$ são dólares por unidade de quantidade; as unidades de $T$ são o número de unidades de quantidade por período de tempo. Podemos converter a equação de uma expressão que se aplica a um \textit{intervalo} de tempo para uma que se aplica a um \textit{ponto} no tempo pelo processo usual de limite, fazendo o intervalo de tempo para o qual agregamos os pagamentos aproximar-se de zero, e expressando $T$ não como um agregado, mas como uma taxa de fluxo (isto é, o limite da razão entre quantidades agregadas e o comprimento do intervalo à medida que o comprimento do intervalo se aproxima de zero). A magnitude $T$ tem então a dimensão de quantidade por unidade de tempo. O produto de $P$ e $T$ tem então a dimensão de dólares por unidade de tempo.

Como o lado direito se destina a resumir um processo contínuo, um fluxo de bens físicos e serviços, o item físico transferido (bem, serviço ou título) é tratado como se desaparecesse da circulação econômica uma vez transferido. Se, por exemplo, um único item, digamos, uma casa, fosse transferido três vezes durante o intervalo de tempo para o qual o $PT$ é medido, ele entraria em $T$ como três casas para esse intervalo de tempo. Além disso, apenas os itens físicos que entram em transações são explicitamente incluídos em $T$. As casas que existem, mas não são compradas ou vendidas durante o intervalo de tempo, são omitidas, embora, se forem alugadas, os valores de aluguel de seus serviços serão incluídos em $PT$ e o número de unidades-aluguel de habitação por ano será incluído em $T$. Claramente, $T$ é um tipo de índice de quantidades bastante especial: ele inclui fluxos de serviços (homens-hora, unidades-ano de habitação, quilowatts-hora), mas também itens de capital que rendem fluxos (casas, usinas elétricas), ponderando cada um desses itens de capital de acordo com o número de vezes que entram em trocas (sua "velocidade de circulação" em estrita analogia com a "velocidade de circulação" da moeda). Da mesma forma, $P$ é um tipo de índice de preços bastante especial.

O lado monetário analisado no lado esquerdo das equações (1) e (2) é tratado de forma muito diferente. A moeda que troca de mãos é tratada como mantendo sua identidade, e toda a moeda, quer usada em transações durante o intervalo de tempo em questão ou não, é explicitamente contabilizada. A moeda é tratada como um estoque, não um fluxo ou uma mistura de um fluxo e um estoque. Para uma única transação, a decomposição em $M$ e $V$ é trivial: o dinheiro que é transferido é girado uma vez, ou $V = 1$. Para todas as transações durante um intervalo, podemos, em princípio, classificar o estoque existente de dólares de moeda de acordo com a entrada de cada dólar em 0, 1, 2, ... transações, isto é, de acordo com quantas vezes cada dólar "girou" 0, 1, 2, ... vezes. A média ponderada desses números de giro, ponderada pelo número de dólares que giraram esse número de vezes, é o conceito conceitual equivalente de $V$. As dimensões de $M$ são dólares; de $V$, número de giros por unidade de tempo; logo, do produto, dólares por unidade de tempo.\footnote{Uma crítica comum à equação quantitativa é que, embora leve em conta a velocidade de circulação da moeda, não leva em conta a velocidade de circulação dos bens. Como deixam claros os dois parágrafos anteriores, embora essa crítica não seja literalmente válida, ela tem um ponto real. A velocidade de circulação da moeda é explícita; a velocidade de circulação dos bens é implícita. Pode muito bem tornar o lado direito das equações (1) e (2) mais significativo torná-lo a soma de dois componentes — um, o valor total das transações envolvendo fluxos contínuos, o outro, o valor das transferências de itens existentes de riqueza — e expressar o segundo componente como um estoque de preços vezes uma velocidade vezes um estoque. Com efeito, a mudança para a versão de renda da equação resolve o problema ao negligenciar completamente as transferências de itens existentes de riqueza.}

A equação (2) difere da equação (1) ao dividir os pagamentos em duas categorias: aqueles efetuados pela transferência de moeda manual (incluindo moedas) e aqueles efetuados pela transferência de depósitos. Na equação (2), $M$ representa apenas o volume de moeda corrente e $V$ a velocidade da moeda corrente, $M'$ o volume de depósitos e $V'$ a velocidade dos depósitos.

Uma razão para a ênfase nesta divisão particular foi a disputa persistente sobre se o termo "moeda" deveria incluir apenas moeda corrente ou também depósitos (Friedman e Schwartz 1970, cap. 2). Outra razão foi a disponibilidade direta de dados sobre $M'V'$ a partir de registros bancários de compensações ou de débitos em contas de depósito. Estes tornam possível calcular $V'$ de uma forma que não é possível calcular $V$.\footnote{Para um cálculo indireto extremamente engenhoso de $V$, não apenas para a moeda como um todo, mas para denominações particulares de moeda corrente, ver Laurent (1969).}

As equações (1) e (2), assim como as outras equações quantitativas que discutirei, destinam-se a ser identidades — uma aplicação especial da contabilidade de partidas dobradas, com cada transação registrada simultaneamente em ambos os lados da equação. No entanto, como ocorre com as identidades da renda nacional com as quais estamos todos familiarizados, quando os dois lados, ou os elementos separados nos dois lados, são estimados a partir de fontes independentes de dados, surgem muitas diferenças entre os dois lados (Mitchell 1927, pp. 128–39). Isso tem sido menos óbvio para as equações quantitativas do que para as identidades da renda nacional — com seu padrão "discrepância estatística" — por causa da dificuldade de calcular $V$ diretamente. Como resultado, $V$ na equação (1) ou $V$ e $V'$ na equação (2) têm sido geralmente calculados como os números que têm a propriedade de tornar as equações corretas. Esses números calculados, portanto, incorporam todo o equivalente à "discrepância estatística".

Assim como o lado esquerdo da equação (1) pode ser dividido em vários componentes, também o lado direito pode sê-lo. A ênfase nas transações refletida nesta versão da equação quantitativa sugere dividir as transações totais em categorias de pagamentos para os quais os períodos ou práticas de pagamento diferem: por exemplo, em transações de capital, compras de bens e serviços finais, compras de bens intermediários, pagamentos pelo uso de recursos, talvez divididos em pagamentos de salários e outros pagamentos. O valor observado de $V$ pode muito bem ser uma função da distribuição dos pagamentos totais entre as categorias. Alternativamente, se a equação quantitativa for interpretada não como uma identidade, mas como uma relação funcional que expressa a velocidade desejada como uma função de outras variáveis, a distribuição dos pagamentos pode muito bem ser um conjunto importante de variáveis.

\subsection{A Forma de Renda da Equação Quantitativa}

Apesar da grande quantidade de trabalho empírico feito sobre as equações de transações, notadamente por Irving Fisher e Carl Snyder (Fisher 1911, pp. 280–318; Fisher 1919; Snyder 1934), as ambiguidades dos conceitos de "transações" e do "nível geral de preços" — particularmente aquelas que surgem da mistura de transações de conta corrente e de capital — nunca foram satisfatoriamente resolvidas. O desenvolvimento mais recente da contabilidade nacional ou social enfatizou as transações de renda em vez das transações brutas e tratou de forma explícita e satisfatória dos problemas conceituais e estatísticos de distinguir entre mudanças nos preços e mudanças nas quantidades. Como resultado, a equação quantitativa tendeu mais recentemente a ser expressa em termos de renda em vez de transações. Seja $Y = $ renda nacional nominal, $P = $ o índice de preços implícito na estimativa da renda nacional a preços constantes, e $y = $ renda nacional a preços constantes, de modo que
\begin{equation}
Y = Py.
\end{equation}

Seja $M$ representando, como antes, o estoque de moeda; mas defina $V$ como o número médio de vezes por unidade de tempo que o estoque de moeda é usado para realizar transações de \textit{renda} (isto é, pagamentos por serviços produtivos finais ou, alternativamente, por bens e serviços finais) em vez de \textit{todas} as transações. Podemos então escrever a equação quantitativa na forma de renda como
\begin{equation}
MV = Py,
\end{equation}
ou, se for desejado distinguir as transações de moeda corrente das de depósito, como
\begin{equation}
MV + M'V' = Py.
\end{equation}

Embora os símbolos $P, V,$ e $V'$ sejam usados tanto nas equações (4) e (5) quanto nas equações (1) e (2), eles representam conceitos diferentes em cada par de equações.

As equações (4) e (5) são conceitualmente e empiricamente mais satisfatórias do que as equações (1) e (2). No entanto, elas têm a desvantagem de negligenciar completamente tanto a proporção de transações intermediárias em relação às finais quanto as transações em ativos de capital existentes.

Na versão de transações da equação quantitativa, cada transação intermediária — isto é, compra de uma empresa por outra — está incluída no valor total da transação, de modo que o valor do trigo, por exemplo, está incluído uma vez quando é vendido pelo agricultor ao moinho, uma segunda vez quando o moinho vende farinha ao padeiro, uma terceira vez quando o padeiro vende pão ao merceeiro, uma quarta vez quando o merceeiro vende pão ao consumidor. Na versão de renda, apenas o valor líquido adicionado por cada uma dessas transações está incluído. Para colocar de outra forma, na versão de transações, o evento elementar é uma troca isolada de um item físico por moeda — um evento real e claramente observável. Na versão de renda, o evento elementar é um evento hipotético que pode ser inferido da observação, mas não é diretamente observável. É uma série completa de transações envolvendo a troca de serviços produtivos por bens finais, através de uma sequência de pagamentos monetários, com todas as transações intermediárias no circuito da renda compensadas. O valor total de todas as transações é, portanto, um múltiplo do valor das transações de renda apenas.

Para um dado fluxo de serviços produtivos ou, alternativamente, de produtos finais (as duas faces da renda), o volume de transações será claramente afetado pela integração ou desintegração vertical das empresas, o que reduz ou aumenta o número de transações envolvidas em um único circuito de renda, ou por mudanças tecnológicas que alongam ou encurtam o processo de transformação de serviços produtivos em produtos finais. O volume de renda não será assim afetado.

Da mesma forma, a versão de transações inclui a compra de um ativo existente — uma casa ou um pedaço de terra ou uma ação ordinária — precisamente em pé de igualdade com uma transação intermediária ou final. A versão de renda exclui tais transações completamente.

Essas diferenças são uma vantagem ou desvantagem da versão de renda? Isso claramente depende do que determina a quantidade de moeda que as pessoas desejam manter. As mudanças do tipo considerado nos parágrafos anteriores, mudanças que alteram a razão entre transações intermediárias e de capital e transações de renda, também alteram na mesma direção e pela mesma proporção a quantidade de moeda que as pessoas desejam manter? Ou elas tendem a deixar essa quantidade inalterada? Ou elas têm um efeito mais complexo?

Claramente, as versões de transações e de renda da teoria quantitativa envolvem concepções muito diferentes do papel da moeda. Para a versão de transações, a coisa mais importante sobre a moeda é que ela é \textit{transferida}. Para a versão de renda, a coisa mais importante é que ela é \textit{mantida}. Essa diferença é ainda mais óbvia na versão de Cambridge de saldos de caixa da equação quantitativa. De fato, a versão de renda pode talvez ser melhor considerada como uma estação intermediária entre a versão de Fisher e a de Cambridge.

\subsection{Abordagem de Saldos de Caixa de Cambridge}

A característica essencial de uma economia monetária é que ela permite que o ato de compra seja separado do ato de venda. Um indivíduo que tem algo a trocar não precisa procurar a dupla coincidência — alguém que tanto queira o que ele tem quanto ofereça em troca o que ele quer. Ele precisa apenas encontrar alguém que queira o que ele tem, vendê-lo a ele por poder de compra geral e, então, encontrar alguém que tenha o que ele quer e comprá-lo com o poder de compra geral.

Para que o ato de compra seja separado do ato de venda, deve haver algo que todos aceitem em troca como "poder de compra geral" — este é o aspecto da moeda enfatizado na abordagem de transações. Mas também deve haver algo que possa servir como uma morada temporária de poder de compra no ínterim entre a venda e a compra. Este é o aspecto da moeda enfatizado na abordagem de saldos de caixa.

Quanto de moeda as pessoas ou empresas desejarão manter para esse propósito? Como primeira aproximação, geralmente se supõe que a quantidade guardada mantém alguma relação com a renda, na suposição de que isso afeta o volume de compras potenciais para as quais o indivíduo ou empresa deseja manter uma morada temporária de poder de compra. Podemos, portanto, escrever
\begin{equation}
M = kPy,
\end{equation}
onde $M, P,$ e $y$ são definidos como na equação (4), e $k$ é a razão entre o estoque de moeda e a renda — seja a razão observada calculada de modo a tornar a equação (6) uma identidade, ou a razão "desejada" de modo que $M$ seja a quantidade "desejada" de moeda, que não precisa ser igual à quantidade real. Em qualquer caso, $k$ é numericamente igual ao recíproco do $V$ na equação (4), sendo o $V$ em um caso interpretado como velocidade medida e no outro como velocidade desejada.

Embora a equação (6) seja simplesmente uma transformação matemática da equação (4), ela realça de forma muito mais nítida a diferença entre os aspectos da moeda enfatizados pela abordagem de transações e aqueles enfatizados pela abordagem de saldos de caixa. Essa diferença faz com que diferentes definições de moeda pareçam naturais e leva a que a ênfase seja colocada em diferentes variáveis e técnicas analíticas.

A abordagem de transações torna natural definir a moeda em termos de tudo o que serve como meio de troca na liquidação de obrigações. Ao enfatizar a função da moeda como uma morada temporária de poder de compra, a abordagem de saldos de caixa torna inteiramente apropriado incluir também tais reservas de valor como depósitos à vista e a prazo não transferíveis por cheque, embora esta abordagem claramente não exija a sua inclusão (Friedman e Schwartz 1970, cap. 3).

Da mesma forma, a abordagem de transações leva a que a ênfase seja colocada em variáveis como práticas de pagamento, arranjos financeiros e econômicos para a realização de transações, e a velocidade de comunicação e transporte, uma vez que afeta o tempo necessário para efetuar um pagamento — essencialmente, para enfatizar os aspectos mecânicos do processo de pagamentos. A abordagem de saldos de caixa, por outro lado, leva a que a ênfase seja colocada em variáveis que afetam a utilidade da moeda como um ativo: os custos e retornos de manter moeda em vez de outros ativos, a incerteza do futuro, e assim por diante — essencialmente, para enfatizar o papel do caixa em uma carteira.

Naturalmente, nenhuma das abordagens impõe a exclusão das variáveis enfatizadas pela outra — e os economistas mais sofisticados que as utilizaram tiveram concepções mais amplas do que a abordagem particular por eles adotada. Os aspectos de carteira entram nos custos de efetuar transações e, consequentemente, afetam os arranjos de pagamento mais eficientes; os aspectos mecânicos entram nos retornos de manter caixa e, portanto, afetam a utilidade do caixa em uma carteira.

Finalmente, com relação às técnicas analíticas, a abordagem de saldos de caixa se ajusta muito mais facilmente ao aparato geral marshalliano de demanda e oferta do que a abordagem de transações. A equação (6) pode ser considerada como uma função de demanda por moeda, com $P$ e $y$ no lado direito sendo duas das variáveis das quais depende a demanda por moeda, e com $k$ simbolizando todas as outras variáveis, de modo que $k$ deve ser considerado não como uma constante numérica, mas como sendo ele próprio uma função de outras variáveis. Para completar, a análise requer outra equação que mostre a oferta de moeda como uma função de outras variáveis. O nível de preços ou o nível da renda nominal é então o resultante da interação das funções de demanda e oferta.

A teoria quantitativa em sua versão de saldos de caixa sugere, assim, organizar a análise dos fenômenos monetários em termos de (1) os fatores que determinam a quantidade nominal de moeda a ser mantida — as condições que determinam a oferta — e (2) os fatores que determinam a quantidade real de moeda que a comunidade deseja manter — as condições que determinam a demanda.

\section{Oferta de Moeda em Unidades Nominais}

Os fatores que determinam a quantidade nominal de moeda disponível para ser mantida dependem criticamente do sistema monetário. Para sistemas como os que prevaleceram nos Estados Unidos e no Reino Unido durante o século passado, eles podem ser utilmente analisados sob os três títulos principais que denominamos os determinantes proximais do estoque de moeda: (1) a quantidade de moeda de alta potência — para qualquer país, esta é determinada através do balanço de pagamentos sob um padrão-mercadoria internacional, pelas autoridades monetárias sob um padrão fiduciário; (2) a razão entre depósitos bancários e as disponibilidades de moeda de alta potência dos bancos — esta é determinada pelo sistema bancário sujeito a quaisquer exigências de reservas impostas a eles por lei ou pelas autoridades monetárias; e (3) a razão entre os depósitos do público e as suas detenções de moeda corrente — esta é determinada pelo público (Friedman e Schwartz 1963b, pp. 776–98; Cagan 1965).

\section{A Demanda por Moeda}

A análise da preferência pela liquidez de J. M. Keynes (discutida adiante na seção 5) reforçou a mudança de ênfase da equação quantitativa para a versão de saldos de caixa — uma mudança de ênfase dos aspectos mecânicos do processo de pagamentos para as qualidades da moeda como um ativo. A análise de Keynes, embora estritamente na tradição de saldos de caixa de Cambridge, foi muito mais explícita ao enfatizar o papel da moeda como um entre muitos ativos e das taxas de juros como o custo relevante de manter moeda.

Trabalhos mais recentes foram ainda mais longe nessa direção, tratando a demanda por moeda como parte da teoria do capital ou da riqueza, preocupada com a composição do balanço ou carteira de ativos.

Deste ponto de vista, é importante distinguir entre detentores finais de riqueza, para quem a moeda é uma forma na qual eles escolhem manter sua riqueza, e empresas, para quem a moeda é um bem de produtor como máquinas ou estoques (Friedman 1956).

\subsection{Demanda por Detentores Finais de Riqueza}

Para os detentores finais de riqueza, a demanda por moeda, em termos reais, pode-se esperar que seja uma função principalmente das seguintes variáveis:

i) \textit{Riqueza total.} — Este é o análogo da restrição orçamentária na teoria usual da escolha do consumidor. É o total que deve ser dividido entre várias formas de ativos. Na prática, estimativas da riqueza total raramente estão disponíveis. Em vez disso, a renda pode servir como um índice de riqueza. No entanto, deve-se reconhecer que a renda conforme medida pelos estatísticos pode ser um índice defeituoso de riqueza porque está sujeita a flutuações erráticas de ano para ano, e um conceito de longo prazo, como o conceito de renda permanente desenvolvido em conexão com a teoria do consumo, pode ser mais útil (Friedman 1957, 1959; Brunner e Meltzer 1963; Meltzer 1963).

A ênfase na renda como um substituto para a riqueza, em vez de como uma medida do "trabalho" a ser feito pela moeda, é conceitualmente talvez a diferença básica entre trabalhos mais recentes e as versões anteriores da teoria quantitativa.

ii) \textit{A divisão da riqueza entre formas humanas e não humanas.} — O principal ativo da maioria dos detentores de riqueza é sua capacidade de ganhos pessoais, mas a conversão de riqueza humana em riqueza não humana ou o inverso está sujeita a limites estreitos por causa de restrições institucionais. Pode-se fazer isso usando ganhos correntes para comprar riqueza não humana ou usando riqueza não humana para financiar a aquisição de habilidades, mas não por compra ou venda e apenas em uma extensão limitada por empréstimo com o colateral dos ganhos futuros. Assim, a fração da riqueza total que está na forma de riqueza não humana pode ser uma variável importante adicional.

iii) \textit{As taxas esperadas de retorno sobre a moeda e outros ativos.} — Este é o análogo dos preços de uma mercadoria e seus substitutos e complementos na teoria usual da demanda do consumidor. A taxa nominal de retorno sobre a moeda pode ser zero, como ocorre geralmente com a moeda corrente, ou negativa, como às vezes ocorre com depósitos à vista sujeitos a tarifas de serviços bancários líquidos, ou positiva, como às vezes ocorre com depósitos à vista sobre os quais juros são pagos e geralmente ocorre com depósitos a prazo. A taxa nominal de retorno sobre outros ativos consiste em duas partes: primeiro, qualquer rendimento ou custo pago ou recebido atualmente, como juros sobre títulos, dividendos sobre ações e custos de armazenamento de ativos físicos, e, segundo, mudanças em seus preços nominais. A segunda parte será, é claro, especialmente importante sob condições de inflação ou deflação.

iv) \textit{Outras variáveis que determinam a utilidade ligada aos serviços prestados pela moeda em relação àqueles prestados por outros ativos} — na terminologia keynesiana, determinando o \textit{valor atribuído à propriedade de liquidez}. — Uma tal variável pode ser uma já considerada — a saber, riqueza real ou renda, visto que os serviços prestados pela moeda podem, em princípio, ser considerados pelos detentores de riqueza como uma "necessidade", como o pão, cujo consumo aumenta menos do que proporcionalmente a qualquer aumento na renda, ou como um "luxo", como o lazer, cujo consumo aumenta mais do que proporcionalmente.

Outra variável que provavelmente será empiricamente importante é o grau de estabilidade econômica esperado para prevalecer no futuro. É provável que os detentores de riqueza atribuam consideravelmente mais valor à liquidez quando esperam que as condições econômicas sejam instáveis do que quando esperam que sejam altamente estáveis. Esta variável é provável que seja difícil de expressar quantitativamente, embora a direção da mudança possa ser clara a partir de informações qualitativas. Por exemplo, o início da guerra claramente produz expectativas de instabilidade, que é uma razão pela qual a guerra é frequentemente acompanhada por um aumento notável nos saldos reais — isto é, um declínio notável na velocidade.

Ainda outra variável pode ser o volume de transferências de capital em relação à renda — do comércio em bens de capital existentes por detentores finais de riqueza. Quanto maior o giro de ativos de capital, maior a fração da riqueza total que as pessoas podem achar útil manter como caixa. Esta variável corresponde à classe de transações negligenciadas na passagem da versão de transações para a de renda da equação quantitativa.

Podemos simbolizar esta análise em termos da seguinte função de demanda por moeda para um detentor individual de riqueza:
\begin{equation}
\frac{M}{P} = f\left( y, w; r_m, r_b, r_e, \frac{1}{P}\frac{dP}{dt}; u \right),
\end{equation}
onde $M, P,$ e $y$ têm o mesmo significado que na equação (6), exceto que se referem a um único detentor de riqueza; $w$ é a fração da riqueza na forma não humana (ou, alternativamente, a fração da renda derivada da propriedade); $r_m$ é a taxa nominal esperada de retorno sobre a moeda; $r_b$ é a taxa nominal esperada de retorno sobre títulos de valor fixo, incluindo mudanças esperadas em seus preços; $r_e$ é a taxa nominal esperada de retorno sobre ações, incluindo mudanças esperadas em seus preços; $(1/P)(dP/dt)$ é a taxa esperada de mudança nos preços de bens e, portanto, a taxa nominal esperada de retorno sobre ativos reais; e $u$ é um símbolo \textit{portmanteau} que representa quaisquer variáveis que não a renda que possam afetar a utilidade ligada aos serviços da moeda. Cada uma das quatro taxas de retorno representa, é claro, um conjunto de taxas de retorno, e para alguns propósitos pode ser importante classificar os ativos de forma ainda mais fina — por exemplo, distinguir a moeda corrente dos depósitos, os títulos de curto prazo dos de longo prazo, os títulos de valor fixo seguros dos arriscados, e um tipo de ativos físicos de outro.\footnote{Sob algumas condições assumidas, as quatro taxas de retorno podem não ser independentes. Por exemplo, em um caso especial considerado em Friedman (1956, pp. 9–10), $r_b = r_e + (1/P)(dP/dt)$.}

Os problemas usuais de agregação surgem ao passar da equação (7) para uma equação correspondente para a economia como um todo — em particular, eles surgem da possibilidade de que a quantidade de moeda demandada possa depender da distribuição entre os indivíduos de variáveis como $y$ e $w$ e não meramente de seu valor agregado ou médio. Se negligenciarmos esses efeitos distributivos, a equação (7) pode ser considerada como aplicando-se à comunidade como um todo, com $M$ e $y$ referindo-se aos saldos monetários e renda real per capita, respectivamente, e $w$ à fração da riqueza agregada na forma não humana.

As principais dificuldades que surgem na prática ao aplicar a equação (7) são as definições precisas de $y$ e $w$, a estimativa das taxas de retorno \textit{esperadas} em contraste com as taxas de retorno reais, e a especificação quantitativa das variáveis designadas por $u$.

\subsection{Demanda por Empresas Comerciais}

As empresas comerciais não estão sujeitas a uma restrição comparável à imposta pela riqueza total dos detentores finais de riqueza. O montante total de capital incorporado em ativos produtivos, incluindo a moeda, é uma variável que pode ser determinada por uma empresa de modo a maximizar os retornos, uma vez que ela pode adquirir capital adicional através do mercado de capitais. Logo, não há razão neste plano para incluir a riqueza total, ou $y$ como um substituto para a riqueza total, como uma variável na função de demanda de moeda das empresas.

Pode, no entanto, ser desejável incluir uma variável um tanto similar definindo a "escala" da empresa em diferentes planos — a saber, como um índice do valor produtivo de diferentes quantidades de moeda para a empresa. Isso é mais próximo da abordagem anterior de transações, enfatizando o "trabalho" a ser feito pela moeda. Não está de forma alguma claro qual é a variável apropriada: transações totais, valor líquido adicionado, renda líquida, capital total em forma não monetária ou patrimônio líquido. A falta de disponibilidade de dados fez com que muito menos trabalho empírico fosse realizado sobre a demanda de moeda das empresas do que sobre uma curva de demanda agregada de moeda que engloba tanto os detentores finais de riqueza quanto as empresas comerciais. Como resultado, há por enquanto apenas indicações vagas sobre a melhor variável a ser usada.

A divisão da riqueza entre as formas humana e não humana não tem relevância especial para as empresas comerciais, uma vez que elas provavelmente comprarão os serviços de ambas as formas no mercado.

As taxas de retorno sobre a moeda e sobre ativos alternativos são, obviamente, altamente relevantes para as empresas comerciais. Essas taxas determinam o custo líquido de manter saldos monetários. No entanto, as taxas particulares que são relevantes podem ser bastante diferentes daquelas que são relevantes para os detentores finais de riqueza. Por exemplo, as taxas cobradas pelos bancos em empréstimos são de menor importância para os detentores de riqueza, mas podem ser extremamente importantes para as empresas, uma vez que os empréstimos bancários podem ser uma forma pela qual elas podem adquirir o capital incorporado nos saldos monetários.

A contraparte para as empresas comerciais da variável $u$ na equação (7) é o conjunto de variáveis que não a escala que afetam a produtividade dos saldos monetários. Pelo menos uma destas — a saber, as expectativas sobre a estabilidade econômica — é provável que seja comum às empresas comerciais e aos detentores finais de riqueza.

Com estas interpretações das variáveis, a equação (7), com $w$ excluído, pode ser considerada como simbolizando a demanda das empresas comerciais por moeda e, como ela está, simbolizando a demanda agregada por moeda, embora com qualificações ainda mais sérias sobre as ambiguidades introduzidas pela agregação.

\section{O Desafio Keynesiano à Teoria Quantitativa}

A análise de renda-despesa desenvolvida por John Maynard Keynes em sua \textit{General Theory} (Keynes 1936) ofereceu uma abordagem alternativa para a interpretação de mudanças na renda nominal que enfatizava a relação entre renda nominal e investimento ou gastos autônomos, em vez da relação entre renda nominal e o estoque de moeda.

O desafio básico de Keynes à teoria reinante pode ser resumido em três proposições que ele estabeleceu:

1. Como uma questão puramente \textit{teórica}, não precisa existir, mesmo que todos os preços sejam flexíveis, uma posição de equilíbrio de longo prazo caracterizada pelo "pleno emprego" de recursos.
2. Como uma questão \textit{empírica}, os preços podem ser considerados como rígidos — um dado institucional — para \textit{flutuações econômicas de curto prazo}; ou seja, para tais flutuações, a distinção entre magnitudes reais e nominais é que o nível de preços é o que está no coração da teoria quantitativa e é de nenhuma importância.
3. A função de demanda por moeda tem uma forma empírica particular — correspondendo à preferência absoluta pela liquidez — que torna a velocidade altamente instável em grande parte do tempo, de modo que mudanças na quantidade de moeda, no mais das vezes, simplesmente produzem mudanças em $V$ na direção oposta. Esta proposição é crítica para ambas as proposições (1) e (2), embora as razões para a preferência absoluta pela liquidez sejam diferentes no longo prazo e no curto prazo. A preferência absoluta pela liquidez a uma taxa de juros que se aproxima de zero é uma condição necessária, embora não suficiente, para a proposição (1). A preferência absoluta pela liquidez à taxa de juros "convencional" explica por que Keynes considerava a equação quantitativa, embora perfeitamente válida como uma identidade, como amplamente inútil para a política ou para prever flutuações de curto prazo na renda nominal e real (idênticas pela proposição [2]). Em seu lugar, Keynes colocou a identidade da renda suplementada por uma propensão estável a consumir.

\subsection{Equilíbrio de Longo Prazo}

A primeira proposição pode ser tratada sumariamente porque se demonstrou ser falsa. O erro de Keynes consistiu em negligenciar o papel da riqueza na função de consumo — ou, dito de outra forma, em negligenciar a existência de um estoque desejado de riqueza como um objetivo que motiva a poupança.\footnote{Keynes, é claro, reconheceu verbalmente este ponto, mas ele não foi incorporado em seu modelo formal da economia. Seu papel fundamental foi apontado primeiro por Haberler (1941, pp. 242, 389, 403, 491–503) e subsequentemente por Pigou (1947), Tobin (1947), Patinkin (1951) e Johnson (1961).}

Todos os tipos de fricções e rigidezes podem interferir na obtenção de uma hipotética posição de equilíbrio de longo prazo em pleno emprego; mudanças dinâmicas na tecnologia, recursos e instituições sociais e econômicas podem mudar continuamente as características dessa posição de equilíbrio; mas não há "falha fundamental no sistema de preços" que torne o desemprego o resultado natural de um mecanismo de mercado plenamente operante.\footnote{Esta proposição desempenhou um papel importante em ganhar para Keynes a adesão de muitos não economistas, particularmente o grande grupo de reformadores, críticos sociais e radicais que estavam convencidos de que havia algo fundamentalmente errado com o sistema "capitalista". Existe uma longa história de tentativas, algumas altamente sofisticadas, para demonstrar que há uma "falha no sistema de preços" (o título de uma dessas tentativas [Martin 1924]), tentativas que remontam pelo menos a Malthus. Em tempos modernos, uma das mais populares e persistentes é a doutrina do "crédito social" do Major C. H. Douglas, que até gerou um partido político no Canadá que tomou o controle (em 1935) do governo de uma das províncias canadenses (Alberta) e tentou implementar alguns dos dogmas de Douglas. Esta política deparou-se com obstáculos legais e teve de ser abandonada. O partido sucessor agora (1969) controla Alberta e a Colúmbia Britânica. Mas, antes de Keynes, essas tentativas tinham sido feitas principalmente por pessoas fora da corrente principal da profissão econômica, e economistas profissionais tinham tido pouca dificuldade em demonstrar suas falhas teóricas e inadequações. A tentativa de Keynes foi, portanto, saudada com entusiasmo. Ela veio de um economista profissional da mais alta reputação, respeitado, e devidamente, por seus colegas economistas como um dos maiores economistas de todos os tempos. O sistema analítico era sofisticado e complexo e, no entanto, uma vez dominado, parecia altamente mecânico e capaz de produzir conclusões de longo alcance e importantes com um mínimo de esforço; e estas conclusões foram, além disso, altamente congeniais aos oponentes do sistema de mercado. Inútil dizer, a demonstração de que esta proposição de Keynes é falsa, e mesmo a aceitação desta demonstração por economistas que se consideravam discípulos de \textit{The General Theory}, não impediu que os oponentes não economistas do sistema de mercado continuassem a acreditar que Keynes provou a proposição e continuassem a citar a sua autoridade para ela.}

\subsection{Rigidez de Preços de Curto Prazo\footnote{Sou devedor a um livro brilhante de Leijonhufvud (1968) para uma plena apreciação da importância desta proposição no sistema de Keynes. Esta subseção e a que segue, sobre a função de preferência pela liquidez, devem muito à análise penetrante de Leijonhufvud.}}

A distinção de Alfred Marshall entre o equilíbrio de mercado, o equilíbrio de curto período e o equilíbrio de longo período era um dispositivo para analisar o ajuste dinâmico em um mercado particular a uma mudança na demanda ou na oferta. Este dispositivo tinha duas características principais. A primeira, a menos importante para os nossos propósitos, é que substituía o processo contínuo por uma série de passos discretos — comparável a aproximar uma função contínua por um conjunto de segmentos de linha reta. A segunda é a suposição de que os preços se ajustam mais rapidamente do que as quantidades, de fato, tão rapidamente que o ajuste de preços pode ser considerado instantâneo. Um aumento na demanda (um deslocamento para a direita da curva de demanda de longo prazo) produzirá um novo equilíbrio de mercado envolvendo um preço mais alto, mas a mesma quantidade. O preço mais alto irá, no curto prazo, encorajar os produtores existentes a produzir mais com suas plantas existentes, aumentando assim a quantidade e trazendo os preços de volta para baixo em direção ao seu nível original, e, no longo prazo, atrair novos produtores e incentivar os produtores existentes a expandir suas plantas, aumentando ainda mais as quantidades e baixando os preços. Ao longo do processo, leva tempo para que o produto se ajuste, mas não leva tempo para os preços o fazerem. Esta suposição não tem efeito sobre a posição de equilíbrio final, mas é vital para o caminho para o equilíbrio.

Esta suposição marshalliana sobre o preço de um produto particular tornou-se amplamente aceita e tendeu a ser transportada impensadamente para o nível de preços na análise do ajuste dinâmico a uma mudança na demanda ou na oferta de moeda. Como observamos acima, a equação de saldos de caixa de Cambridge presta-se a uma interpretação de demanda-oferta ao longo de linhas marshallianas (Pigou 1917). Assim interpretada, uma mudança na quantidade nominal de moeda (um deslocamento de uma só vez na escala de oferta) exigirá uma mudança em uma ou mais das variáveis do lado direito da equação (6) — $k$ ou $P$ ou $y$ — a fim de reconciliar a demanda e a oferta. No equilíbrio final pleno, o ajuste irá, em geral, residir inteiramente em $P$, uma vez que a mudança na quantidade nominal de moeda não precisa alterar nenhum dos fatores "reais" dos quais $k$ e $y$ em última instância dependem.\footnote{O "em geral" é inserido para alertar o leitor de que esta é uma questão complexa, exigindo para uma análise completa um enunciado muito mais cuidadoso de exatamente como a quantidade de moeda é aumentada. No entanto, estas questões mais sofisticadas não são relevantes para o ponto sob discussão e, portanto, são contornadas.} Na caixa marshalliana, a posição final não é afetada pelas velocidades relativas de ajuste.

Não há nada na lógica da teoria quantitativa que especifique o caminho dinâmico do ajuste, nada que exija que o ajuste total ocorra através de $P$ em vez de através de $k$ ou $y$. Foi amplamente reconhecido que o ajuste durante o que Fisher, por exemplo, chamou de "períodos de transição" ocorreria na prática parcialmente em $k$ e em $y$, bem como em $P$. No entanto, este reconhecimento não foi incorporado na análise teórica formal. A análise formal simplesmente assumiu a suposição de Marshall. Neste sentido, os teóricos quantitativos podem ser validamente criticados por terem "assumido" a flexibilidade de preços — assim como Keynes pode ser validamente criticado por "assumir" que o consumo é independente da riqueza, embora ele tenha reconhecido em seus comentários que a riqueza tem um efeito sobre o consumo.

Keynes foi um verdadeiro marshalliano no método. Ele seguiu Marshall ao tomar a análise de demanda-oferta como sua estrutura. Ele seguiu Marshall ao substituir o ajuste contínuo por uma série de passos discretos e, assim, analisar um processo dinâmico em termos de uma série de deslocamentos entre posições de equilíbrio estático. Mesmo seus passos eram essencialmente de Marshall, sendo sua distinção de curto prazo baseada na fixidez do estoque de capital agregado. No entanto, ele tendeu a fundir o período de mercado e o período de curto prazo e, fiel ao seu próprio dito enganoso, "no longo prazo estamos todos mortos", ele concentrou-se quase exclusivamente no curto prazo.

Keynes também seguiu Marshall ao assumir que uma variável se ajustava tão rapidamente que o ajuste poderia ser considerado instantâneo, enquanto a outra variável se ajustava lentamente. Onde ele se desviou de Marshall, e foi um desvio momentâneo, foi em inverter os papéis atribuídos ao preço e à quantidade. Ele assumiu que, pelo menos para mudanças na demanda agregada, a quantidade era a variável que se ajustava rapidamente, enquanto o preço era a variável que se ajustava lentamente,\footnote{Referi-me a "quantidade", não a "produto", porque conjeturo que Keynes, se pressionado a distinguir o mercado do período de curto prazo, teria feito isso considerando a quantidade disponível para compra como ajustando-se rapidamente no período de mercado em grande parte através de mudanças nos estoques, e no período de curto prazo através de mudanças no produto.} pelo menos em uma direção descendente. Keynes incorporou esta suposição em seu modelo formal expressando todas as variáveis em unidades de salário, de modo que sua análise formal — à parte algumas referências passageiras a uma situação de inflação "verdadeira" — lidava com magnitudes "reais", não magnitudes "nominais" (Keynes 1936, pp. 119, 301, 303). Ele racionalizou a suposição em termos de rigidez salarial decorrente em parte da ilusão monetária, em parte da força dos sindicatos. E, a um nível ainda mais profundo, ele racionalizou a rigidez salarial pela proposição (1): sob condições em que não havia equilíbrio de pleno emprego, não havia nenhum nível de preço nominal de equilíbrio; algo tinha que ser trazido de fora para fixar o nível de preços; poderia muito bem ser a rigidez salarial institucional. Dito de outra forma, os salários nominais flexíveis sob tais circunstâncias não tinham nenhuma função econômica a desempenhar; logo, eles poderiam muito bem ser rígidos.

No entanto, racionalizada, a razão básica para a suposição foi indubitavelmente a falta de concordância entre os fenômenos observados e as implicações de uma aplicação literal da suposição de Marshall às magnitudes agregadas. Tal aplicação literal implicava que as flutuações econômicas tomariam a forma inteiramente de flutuações nos preços com pleno emprego contínuo de homens e recursos. Claramente, isso não correspondia à experiência. Se algo, pelo menos na década e meia entre o fim da Primeira Guerra Mundial e a redação de \textit{The General Theory}, as flutuações econômicas manifestavam-se em maior grau no produto e no emprego do que nos preços. Parecia, portanto, altamente plausível que, pelo menos para fenômenos agregados, as velocidades relativas de ajuste fossem justamente o inverso daquelas assumidas por Marshall.\footnote{Não quero sugerir que a suposição de Marshall seja sempre a melhor para mercados particulares. Pelo contrário, um dos avanços mais sofisticados nos últimos anos na teoria dos preços relativos é o desenvolvimento de modelos de ajuste de preços mais sofisticados que permitem que as taxas de ajuste tanto do preço quanto da quantidade variem continuamente entre o ajuste instantâneo e o muito lento. No entanto, esses desenvolvimentos não são diretamente relevantes para a presente discussão, embora inspirem parcialmente a seção 8 adiante.}

Keynes explorou esta percepção penetrante levando-a ao extremo: todo o ajuste na quantidade, nenhum no preço. Ele qualificou este enunciado assumindo que ele se aplicava apenas a condições de subemprego. No "pleno" emprego, ele mudou para a teoria quantitativa e afirmou que todo ajuste seria no preço — ele designou esta situação de inflação "verdadeira". No entanto, Keynes não deu mais do que um serviço da boca para fora a esta possibilidade, e ele e seus discípulos fizeram o mesmo; por isso, não desvirtua o corpo de sua análise negligenciar amplamente a qualificação.

Dada esta suposição, uma mudança na quantidade nominal de moeda significa uma mudança na quantidade real de moeda. Na equação (6), podemos dividir por $P$, tornando o lado esquerdo a quantidade real de moeda. Uma mudança em $M$ (nominal e real) na quantidade de moeda terá então de ser compensada por uma mudança em $k$ ou em $y$.

Nada até este ponto parece impedir Keynes de ter uma teoria puramente monetária das flutuações econômicas, com mudanças em $M$ sendo refletidas inteiramente em $y$. No entanto, isso entrava em conflito com a interpretação de Keynes dos fatos da Grande Depressão, que ele considerava, acredito erroneamente, como mostrando que a política monetária expansionista era ineficaz em deter um declínio (Friedman 1967). Daí, ele estava inclinado a interpretar mudanças em $M$ como sendo refletidas em $k$ em vez de em $y$. É aqui que sua proposição (3) sobre a preferência pela liquidez entra.

De fato, na forma mais extrema e, sou tentado a dizer, mais pura de sua análise, Keynes supõe que todo o ajuste será em $k$. E, curiosamente, este resultado também pode ser considerado como uma consequência direta de sua suposição sobre a velocidade relativa de ajuste de preço e quantidade. Pois $k$ não é uma constante numérica, mas uma função de outras variáveis. Ele incorpora a preferência pela liquidez. No sistema de Keynes, a variável principal de que depende é a taxa de juros. Esta é um preço. Portanto, era natural para Keynes considerá-la lenta para ajustar, e tomar, como a variável que responde, a quantidade real de moeda que as pessoas desejam manter.

Se mudanças em $M$ não produzem mudanças em $y$, o que produz? A resposta de Keynes é a necessidade de reconciliar o montante que algumas pessoas querem gastar para adicionar ao estoque de capital produtivo com o montante que a comunidade quer poupar para adicionar ao seu estoque de riqueza. Daí Keynes coloca no centro de sua análise a distinção entre consumo e poupança, ou mais fundamentalmente, entre gastos ligados de perto à renda corrente e gastos que são amplamente independentes da renda corrente.

Como resultado tanto da experiência quanto de uma análise teórica adicional, hoje dificilmente há um economista que aceite a conclusão de Keynes sobre o caráter estritamente passivo de $k$, ou a conclusão acompanhante de que a moeda (no sentido da quantidade de moeda) não importa, ou que afirmará explicitamente que $P$ é "realmente" um dado institucional que será completamente inafetado em curtos períodos por mudanças em $M$ (Friedman 1968).

No entanto, a suposição de Keynes sobre a velocidade relativa de ajuste do preço e da quantidade ainda é a chave para a diferença na abordagem e na análise entre aqueles economistas que se consideram keynesianos e aqueles que não se consideram. Independentemente do que o primeiro grupo possa dizer em suas notas laterais e em suas qualificações, eles tratam o nível de preços como um dado institucional em sua análise teórica formal. Eles continuam a considerar mudanças na quantidade nominal de moeda como equivalentes a mudanças na quantidade real de moeda e, portanto, como tendo de ser refletidas em $k$ e $y$. E continuam a considerar o efeito inicial como sendo sobre $k$. A diferença é que eles não consideram mais as taxas de juros como dados institucionais, como Keynes fez em medida considerável. Em vez disso, eles consideram a mudança em $k$ como exigindo uma mudança nas taxas de juros que, por sua vez, produz uma mudança em $y$. Daí, eles atribuem mais significado a mudanças na quantidade de moeda do que Keynes e seus discípulos o fizeram na primeira década ou pouco mais após o surgimento de \textit{The General Theory}.

Uma ilustração notável é fornecida em uma monografia recente da Fundação Cowles, editada por Donald Hester e James Tobin, sobre \textit{Financial Markets and Economic Activity} (Hester e Tobin 1967). Um ensaio fundamental naquele livro apresenta uma análise estática comparativa do ajuste de equilíbrio geral dos estoques de ativos. No entanto, a distinção entre magnitudes nominais e reais nem sequer é discutida. Toda a análise é válida apenas sob a suposição implícita de que os preços nominais de bens e serviços são completamente rígidos, embora as taxas de juros e as magnitudes reais sejam flexíveis.\footnote{Ver Tobin e Brainard (1967). Um exemplo específico que documenta esta afirmação é que Tobin e Brainard assumem explicitamente que os bancos centrais podem determinar a razão entre a moeda corrente (ou moeda de alta potência) e a riqueza total, incluindo ativos reais (Hester e Tobin 1967, pp. 61–62). Se os preços forem flexíveis, o banco central pode determinar apenas magnitudes nominais, não uma tal razão real. Outros artigos na Monografia 21, notadamente o artigo de Brainard, "Financial Institutions and a Theory of Monetary Control" (Brainard 1967), fazem as mesmas suposições implícitas. A palavra "preços" não aparece no índice cumulativo de assuntos desta monografia e de dois volumes companheiros, Monografias 19 e 20.}

Ainda outro exemplo mais recente é um artigo dos mesmos autores, "Pitfalls in Financial Model Building" (Tobin e Brainard 1968), no qual eles apresentam uma simulação de uma "economia fictícia de nossa construção". Nesta economia, o valor de reposição dos ativos físicos é usado como o numerário do sistema, e todos os preços são expressos em relação ao valor de reposição. O resultado é que o sistema — destinado a iluminar os problemas da análise monetária — toma o nível de preços absoluto como determinado fora do sistema. O Banco Central é implicitamente assumido como capaz de determinar o volume \textit{real} e não meramente o nominal das reservas bancárias.

Outro exemplo notável é Gramley e Chase (1965). Neste artigo, a suposição sobre a rigidez de preços é explícita e apresentada como se fosse apenas uma suposição tentativa feita para conveniência da análise. No entanto, o significado empírico que Gramley e Chase atribuem aos seus resultados desmente esta profissão.

Veja também o estudo econométrico de Goldfeld (1966), que se concentra nas formas reais das funções estimadas por causa da "superioridade da versão deflacionada" (p. 166).

A evidência para um período um tanto anterior é fornecida por Holzman e Bronfenbrenner (1963). As teorias de inflação decorrentes da abordagem keynesiana enfatizam fatores institucionais, e não monetários. Elas não teriam aparecido como questões em aberto, e grandes partes delas nunca teriam sido escritas, tivéssemos nós, implícita ou explicitamente, aceitado a suposição de Keynes de que os preços são um dado institucional.

\subsection{Preferência Absoluta pela Liquidez}

Keynes deu uma forma altamente específica à equação (6) ou (7). A quantidade de moeda demandada, ele argumentou, poderia ser tratada como se fosse dividida em duas partes, uma parte, $M_1$, "mantida para satisfazer os motivos transacional e de precaução", a outra, $M_2$, "mantida para satisfazer o motivo especulativo" (Keynes 1936, p. 199). Ele considerava $M_1$ como uma fração aproximadamente constante da renda. Considerava o (curto prazo) da demanda por $M_2$ como surgindo da "incerteza quanto ao futuro da taxa de juros" e a quantidade demandada como dependendo da relação entre as taxas de juros atuais e as taxas de juros esperadas para prevalecer no futuro (Keynes 1936, p. 168; itálicos no original). Keynes, é claro, enfatizou que havia todo um complexo de taxas de juros. No entanto, por simplicidade, ele falou em termos de "a taxa de juros", geralmente significando com isso a taxa sobre títulos de longo prazo que envolviam riscos mínimos de inadimplência — por exemplo, títulos do governo. A distinção fundamental para Keynes era entre títulos de curto prazo e de longo prazo, não entre títulos fixados em valor nominal e aqueles que não eram. Esta última distinção foi tornada irrelevante pela sua suposição de que os preços eram rígidos.

A distinção entre títulos de curto prazo e de longo prazo era importante para Keynes porque correspondia a diferenças no risco de ganho ou perda de capital como resultado de mudanças nas taxas de juros. Para títulos de curto prazo, mudanças nas taxas de juros teriam pouco efeito. Para títulos de longo prazo, o efeito é importante. Leijonhufvud argumentou, e acredito corretamente, que Keynes usou o termo "moeda" referindo-se não apenas à moeda corrente e depósitos estritamente definidos, mas a toda a gama de ativos de curto prazo que forneciam "liquidez" no sentido de segurança contra perda de capital decorrente de mudanças nas taxas de juros.\footnote{Escusado será dizer que Keynes também considerava outros tipos de riscos, como riscos de inadimplência, como altamente relevantes, mas, consistente com a sua proposição (2), ele desconsiderou quase inteiramente os riscos decorrentes de mudanças no nível de preços de bens e serviços (Leijonhufvud 1968, cap. 2).}

É, portanto, um tanto enganoso considerar Keynes, como a maioria da literatura faz, como distinguindo entre "moeda" e "títulos". No entanto, continuarei a seguir a prática corrente e usarei essa terminologia. Uma justificativa para fazê-lo é que Keynes de fato tratou os ativos de curto prazo que rotulou como "moeda" como rendendo retorno de juros zero.\footnote{Neste aspecto, o Comitê Radcliffe é fiel a Keynes ao tratar a "liquidez" amplamente definida como o agregado monetário relevante, em vez da "moeda" estritamente definida.} (É bom recordar que ele estava escrevendo em um momento em que as taxas de juros de curto prazo eram extremamente baixas tanto absolutamente quanto relativamente às taxas de longo prazo. Seu procedimento pareceria altamente irrealista hoje.)

Para formalizar a análise de Keynes nos termos dos símbolos que usamos até agora, podemos escrever sua função de demanda como
\begin{equation}
\frac{M}{P} = \frac{M_1}{P} + \frac{M_2}{P} = k_1 y + f(i - i^*, i^*),
\end{equation}
onde $i$ é a taxa de juros atual, $i^*$ é a taxa de juros que se espera prevalecer, e $k_1$, o análogo do recíproco da velocidade-renda de circulação da moeda, é tratado como determinado pelas práticas de pagamento e, portanto, como uma constante pelo menos no curto prazo.\footnote{Escritores posteriores nesta tradição argumentaram que $k_1$ também deve ser considerado como uma função das taxas de juros. Ver Baumol (1952) e Tobin (1956). No entanto, esta questão não é relevante para a presente discussão.} A taxa de juros atual, $i$, é uma magnitude observada. Logo, ela será a mesma para todos os detentores de moeda, se, como Keynes, abstrairmos da existência de um complexo de taxas de juros. A taxa esperada, $i^*$, não é observável. Ela pode diferir de um detentor para outro e, para cada detentor separadamente, deve ser interpretada como o valor médio de uma distribuição de probabilidade, não como um único valor antecipado com certeza. Para uma função agregada, $i^*$ deve rigorosamente ser interpretada como um vetor, não como um número. Embora eu tenha introduzido $P$ na equação para consistência com minhas equações anteriores, Keynes o omitiu porque sua proposição (2) significava que $P$ era tomado como uma constante.

Em um "dado estado de expectativas", isto é, para um dado valor de $i^*$, quanto maior for a taxa de juros atual, menor será a quantidade de moeda que as pessoas desejarão manter por motivos especulativos. O custo de manter moeda em vez de títulos seria maior de duas maneiras: primeiro, um montante maior de ganhos correntes seria sacrificado; segundo, seria mais provável que as taxas de juros caíssem e, portanto, os preços dos títulos subissem, e assim um montante maior de ganhos de capital seria sacrificado.

Embora as expectativas recebam grande destaque no desenvolvimento da função de preferência pela liquidez que expressa a demanda por $M_2$, Keynes e seus seguidores geralmente não as introduziram explicitamente, como eu fiz, naquela função. No mais das vezes, Keynes e seus seguidores na prática trataram o montante de $M_2$ demandado simplesmente como uma função da taxa de juros atual, a ênfase nas expectativas servindo apenas como uma razão para sua atribuição de instabilidade à função de liquidez.\footnote{Uma exceção notável é Tobin (1958, pp. 65–86).}

A razão para esta omissão é sua concentração na função de demanda de curto prazo. Para essa função, eles consideravam $i^*$ como fixo, de modo que a demanda especulativa era uma função apenas de $i$. Eu introduzi $i^*$ a fim de distinguir entre as diferentes razões que estão implícitas na análise de Keynes para a preferência absoluta pela liquidez no curto prazo e no longo prazo.

O toque especial de Keynes foi menos expressar a função de demanda na forma geral descrita pela equação (8) do que a forma particular que ele deu à função $f(i - i^*, i^*)$. Para um dado $i^*$, ele acreditava que esta função seria altamente elástica em $i = i^*$, sendo o grau de elasticidade em um valor numérico observado de $i$ dependente de quão homogêneas as expectativas de diferentes detentores de moeda são e quão firmemente elas são mantidas.\footnote{Tobin (1958) apresenta uma análise excelente e iluminadora deste caso. Como ele assume que os deslocamentos para dentro ou para fora de títulos envolvem compromissos por um período finito igual à unidade de tempo em termos de que a taxa de juros é expressa, seu valor crítico não é $i = i^*$, mas $i = i^*/(1 + i^*)$, sendo a renda corrente sobre os títulos compensada por uma perda de capital esperada.} Seja um subgrupo substancial de detentores de moeda que têm a mesma expectativa e que mantêm essa expectativa firmemente, e $f$ se tornará perfeitamente elástica àquela taxa de juros atual. Moeda e títulos se tornariam substitutos perfeitos; a preferência pela liquidez se tornaria absoluta. As autoridades monetárias achariam impossível mudar a taxa de juros porque os especuladores detendo essas expectativas firmes as frustrariam.

Sob tais circunstâncias, se as autoridades monetárias buscassem aumentar a quantidade de moeda comprando títulos, isso tenderia a elevar os preços dos títulos e baixar a taxa de retorno sobre os títulos. Mesmo o menor abaixamento, argumentou Keynes, levaria os especuladores com expectativas firmes a absorver os saldos monetários adicionais e vender os títulos demandados pelos detentores de moeda. O resultado seria simplesmente que a comunidade como um todo estaria disposta a manter a quantidade aumentada de moeda; $k$ seria maior e $V$ menor. Inversamente, se as autoridades monetárias diminuíssem a quantidade de moeda vendendo títulos, isso tenderia a elevar a taxa de juros, e mesmo o menor aumento induziria os especuladores a absorver os títulos oferecidos. (Na análise de Keynes, o resultado seria o mesmo se a quantidade de moeda fosse aumentada ou diminuída por operações que adicionassem ou subtraíssem da riqueza total, em vez de substituir uma forma de riqueza por outra, porque ele assumiu que a riqueza não tinha efeito direto nos gastos.)

Ou, novamente, suponha que haja um aumento na renda nominal por qualquer razão. Isso exigirá um aumento em $M_1$, que pode vir de $M_2$ sem quaisquer efeitos adicionais. Inversamente, qualquer declínio em $M_1$ pode ser adicionado a $M_2$ sem quaisquer efeitos adicionais. A conclusão é que sob circunstâncias de preferência absoluta pela liquidez, a renda pode mudar sem uma mudança em $M$ ou nas taxas de juros e $M$ pode mudar sem uma mudança na renda ou nas taxas de juros. Os detentores de moeda estão em equilíbrio metaestável, como um copo de lado sobre uma superfície plana; eles ficarão satisfeitos com qualquer montante de moeda que por acaso possuam.

Para a escala de demanda de longo prazo, a razão para a preferência pela liquidez é diferente. No equilíbrio de longo prazo, $i$ deve ser igual a $i^*$, de modo que $f(i - i^*, i^*)$ reduz-se a uma função apenas de $i^*$. Seja uma deficiência de oportunidades de investimento, o tipo de situação previsto na proposição (1) de Keynes, de modo que $i^*$ se torne muito baixo. Quanto menor a taxa, menor o retorno dos ativos de capital que não a moeda — sejam eles títulos, ações ou ativos físicos (recorde-se que por causa da suposição de que o nível de preços é rígido, Keynes não considerava a distinção entre estes como importante). Consequentemente, quanto menor o $i^*$, menor o custo de manter moeda. Em um nível suficientemente baixo, mas ainda positivo, o retorno extra de manter ativos não monetários apenas compensaria os riscos extras envolvidos. Logo, a essa taxa, a preferência pela liquidez seria absoluta. A "taxa de mercado" de juros não poderia cair indefinidamente; um limite inferior foi estabelecido pelo desejo generalizado de substituir outros ativos por moeda a taxas de juros baixas.

Esta conclusão foi um elemento-chave na proposição (1) de Keynes. Uma maneira de resumir seu argumento para essa proposição é em termos de um possível conflito entre a taxa de juros "de mercado" e a "de equilíbrio". Se as oportunidades de investimento fossem escassas, mas o desejo do público de poupar fosse forte, a taxa de juros de "equilíbrio", ele argumentou, poderia ter de ser muito baixa ou mesmo negativa para igualar o investimento e a poupança. Mas havia um piso para a "taxa de mercado" estabelecido pela preferência pela liquidez. Se este piso excedesse a "taxa de equilíbrio", argumentou ele, havia um conflito que só poderia ser resolvido pelo desemprego que frustrasse a parcimônia do público. A falácia neste argumento é que a introdução da moeda não apenas introduz um piso para a "taxa de mercado"; ela também estabelece um piso para a "taxa de equilíbrio". E, no longo prazo, os dois pisos são idênticos. Esta é a essência do chamado efeito Pigou (Friedman 1962, pp. 262–63).

Nem o próprio Keynes, nem a maioria de seus discípulos e seguidores, distinguiram tão nitidamente quanto eu fiz entre armadilhas de liquidez de curto prazo e de longo prazo. Eles tenderam a fundir as duas e, em linha com a ênfase geral no curto prazo, a enfatizar a elasticidade em relação à taxa de juros atual, não esperada.\footnote{Tobin faz uma distinção explícita deste tipo, embora não em conexão com uma armadilha de liquidez propriamente dita.}

Keynes considerava a preferência absoluta pela liquidez como um caso estritamente "limitante" que, embora "possa tornar-se praticamente importante no futuro", ele sabia de "nenhum exemplo... até agora" (Keynes 1936, p. 207). No entanto, ele tratou a velocidade como se na prática o seu comportamento frequentemente se aproximasse do que prevaleceria neste caso limitante.

A preferência absoluta pela liquidez não é mais explicitamente defendida pelos economistas de hoje — o fracasso dos bancos centrais em suas tentativas de fixar taxas de juros a níveis baixos tornou essa proposição insustentável. No entanto, como os preços absolutamente rígidos, ela ainda desempenha um papel importante na teorização de muitos economistas. Está implícita na tendência de considerar $k$ ou a velocidade como se ajustando passivamente a mudanças na quantidade de moeda. É explícita na tendência de considerar a demanda por moeda como "altamente" elástica em relação às taxas de juros.

Considere novamente a equação (6). Seja uma mudança em $M$. Os economistas na tradição keynesiana continuam, como observado anteriormente, a considerar $P$ como um dado institucional e assim inafetado. Eles devem, portanto, considerar a mudança em $M$ como afetando $k$ ou $y$ ou ambos. Com a preferência absoluta pela liquidez, $k$ pode absorver o impacto sem qualquer mudança na taxa de juros. Uma vez que eles tomam a taxa de juros como o único elo entre a mudança monetária e a renda real, toda a mudança seria então absorvida em $k$ sem efeito sobre $y$. Se a preferência pela liquidez não for absoluta, $k$ pode mudar apenas através de uma mudança na taxa de juros. Mas isso tem efeitos sobre $y$ através dos gastos de investimento. Quanto mais elástica a demanda por moeda, menor a mudança nas taxas de juros terá de ser. Quanto mais inelásticos os gastos de investimento e a poupança em relação à taxa de juros, menor será o efeito de qualquer mudança na taxa de juros sobre $y$. Logo, a tendência de os economistas considerarem $k$ como absorvendo o impacto principal das mudanças em $M$ significa que eles implicitamente ou explicitamente consideram a demanda por moeda como altamente elástica em relação à taxa de juros e os gastos de investimento e poupança como altamente inelásticos.

A tendência por parte de muitos economistas de assumir implicitamente que os preços são um dado institucional e que a demanda por moeda é altamente elástica em relação à taxa de juros subjaz a algumas das críticas que foram dirigidas contra trabalhos anteriores meus e de associados. Fomos interpretados, erroneamente, creio eu, como dizendo que $k$ é completamente independente das taxas de juros (Friedman 1966). Nesse caso, as mudanças em $M$ precisam ser refletidas em $k$ se, também, $P$ é tomado como um dado institucional, todo o efeito será sobre $y$. Esta é a fonte implícita da crítica dirigida contra nós, de que consideramos a quantidade de moeda como determinante do nível de atividade econômica. Não apenas, dizem nossos críticos, acreditamos que a moeda importa, acreditamos que a moeda é tudo o que importa (Okun 1963; Tobin 1965a, p. 481).

Se $P$ não for considerado como um dado institucional, e não o consideramos assim, mesmo se supuséssemos que $k$ fosse completamente insensível às taxas de juros e a qualquer outra coisa que pudesse ser afetada por mudanças em $M$ (como a taxa de mudança em $P$ ou em $y$) e assim fosse uma constante absoluta, à parte distúrbios aleatórios, algo que não a quantidade de moeda teria de ser trazido para a análise para explicar quanto da mudança em $M$ seria refletida em $P$ e quanto em $y$ (veja seção 8, abaixo).

Sempre tentamos qualificar nossas afirmações sobre a importância das mudanças em $M$ referindo-nos ao seu efeito na renda nominal. Mas esta qualificação pareceu sem sentido para os economistas que implicitamente identificaram o nominal com magnitudes reais. Logo, eles entenderam mal nossas conclusões.

Aceitamos a presunção da teoria quantitativa, e pensamos que ela é apoiada pela evidência que examinamos, de que mudanças na quantidade de moeda como tais no longo prazo têm um efeito negligenciável sobre a renda real, de modo que forças não monetárias são "tudo o que importa" para mudanças na renda real ao longo das décadas e a moeda "não importa". Por outro lado, consideramos a quantidade de moeda, mais as outras variáveis (incluindo a própria renda real) que afetam $k$ como sendo essencialmente "tudo o que importa" para a determinação de longo prazo da renda nominal. O preço de mercado é então um resultado conjunto das forças monetárias que determinam a renda nominal e das forças reais que determinam a renda real (Friedman 1958, pp. 242–46; Friedman e Schwartz 1963b, p. 695).

Para períodos de tempo mais curtos, argumentamos que mudanças em $M$ serão refletidas em todas as três variáveis no lado direito da equação (6): $k, P,$ e $y$. Mas argumentamos que o efeito sobre $k$ é empiricamente não para absorver a mudança em $M$, como a análise keynesiana implica, mas frequentemente para reforçá-la, mudanças em $M$ e $k$ frequentemente afetando a renda na mesma direção em vez de direções opostas. Logo, enfatizamos que mudanças em $M$ são um fator principal, embora mesmo assim não o único fator, que explica flutuações de curto prazo tanto na renda nominal quanto no nível real de atividade ($y$). Considero a descrição de nossa posição como "a moeda é tudo o que importa para mudanças na renda nominal e para flutuações de curto prazo na renda real" como um exagero, mas que dá o sabor correto de nossas conclusões. Considero a afirmação de que "a moeda é tudo o que importa", ponto final, como uma deturpação básica de nossas conclusões (Friedman 1958, pp. 246–51; Friedman e Schwartz 1963a, pp. 38–39, 45–46, 55–64; Friedman e Schwartz 1963b, p. 678).

\section{Um Modelo Comum Simples}

Podemos resumir os pontos principais das seções precedentes deste artigo e lançar as bases para as seções finais, estabelecendo um modelo agregado altamente simplificado de uma economia que engloba tanto uma teoria quantitativa simplificada quanto uma teoria simplificada de renda-despesa como casos especiais.

Para este propósito, podemos negligenciar o comércio exterior, assumindo uma economia fechada, e o papel fiscal do governo, assumindo que não há nem gastos governamentais nem receitas governamentais. Também podemos negligenciar distúrbios estocásticos. No que me concentrarei é na divisão da renda nacional entre gastos induzidos e autônomos e no ajuste entre a demanda e a oferta de moeda.

O modelo simples é dado por seis equações:
\begin{equation}
\frac{C}{P} = f\left(\frac{Y}{P}, r\right);
\end{equation}
\begin{equation}
\frac{I}{P} = g(r);
\end{equation}
\begin{equation}
\frac{Y}{P} = \frac{C}{P} + \frac{I}{P} \left(\text{ou, alternativamente, } \frac{S}{P} = \frac{Y - C}{P} = \frac{I}{P}\right);
\end{equation}
\begin{equation}
M^D = P \cdot l\left(\frac{Y}{P}, r\right);
\end{equation}
\begin{equation}
M^S = h(r);
\end{equation}
\begin{equation}
M^D = M^S.
\end{equation}

As três primeiras equações descrevem o ajuste dos fluxos de poupança e investimento; as três últimas, do estoque de moeda demandada e ofertada. A equação (9) é uma função consumo (a "propensão marginal a consumir" de Keynes) expressando o consumo real ($C/P$) como uma função da renda real ($Y/P = y$) e da taxa de juros ($r$). Por simplicidade, a riqueza é omitida, embora, se o modelo fosse usado para ilustrar a proposição (1) de Keynes, e por que ela é falaciosa, a riqueza teria de ser incluída como um argumento na função.

A equação (10) é uma função investimento (a eficiência marginal do investimento de Keynes) que expressa o investimento real ($I/P$) como uma função da taxa de juros. Aqui novamente, consistente com Keynes e a literatura subsequente, tanto o estoque total de capital quanto a renda real poderiam ser incluídos como argumentos. No entanto, no espírito de Keynes, o modelo refere-se a um período curto no qual o estoque de capital pode ser considerado fixo. Para um modelo de período mais longo, o estoque de capital teria de ser incluído e tratado como uma variável endógena, presumivelmente definida por uma integral de investimentos passados. A inclusão da renda na equação, como uma variável independente, confundiria o ponto chave da distinção entre $C$ e $I$. Como matéria teórica, a distinção relevante não é entre consumo e investimento, mas entre gastos que estão estreitamente ligados à renda corrente ("condicionais" na renda seriam, deste ponto de vista, um mnemônico melhor para $C$ do que consumo, embora o termo usualmente usado seja "induzidos") e gastos que são autônomos, isto é, independentes (um mnemônico melhor para $I$ do que investimento) da renda. A identificação dessas categorias com consumo e investimento é uma hipótese empírica. Para propósitos teóricos, qualquer parte do gasto de investimento que seja condicional na renda corrente deveria ser incluída em $C$.

A equação (11) é tipicamente referida como a identidade da renda. Como a transformação entre parênteses deixa claro, ela também pode ser considerada como uma equação de ajuste ou de equilíbrio de mercado especificando que a poupança deve ser igual ao investimento.

A equação (12) é a função de demanda por saldos monetários nominais (a função de preferência pela liquidez de Keynes). É simplesmente a equação (6) ou (7) reescrita em forma simplificada e expressa a quantidade real de moeda demandada ($M^D/P$) como uma função da renda real e da taxa de juros. Aqui novamente, como na equação (9), a riqueza poderia ser propriamente incluída, mas é omitida por simplicidade.

A equação (13) é a função de oferta de moeda nominal. Para ser consistente com a literatura, a taxa de juros entra como uma variável. No entanto, para o que se segue, usaremos o modelo sem ser afetado de nenhuma forma por tratar $M^S$ simplesmente como uma variável exógena, determinada, digamos, pelas autoridades monetárias.\footnote{Isso seria consistente com as descobertas de Cagan sobre a ausência de qualquer efeito significativo das mudanças na taxa de juros sobre a oferta de moeda. No entanto, para ser consistente com suas descobertas, a renda ou algum outro indicador de ciclos econômicos teria de ser incluído como uma variável, como foi feito em alguns estudos empíricos da oferta de moeda. Ver Cagan (1965, pp. 150, 228–32) e Hendershott (1968).}

A equação (14) é a contraparte da equação (11), uma equação de ajuste ou de equilíbrio de mercado especificando que a moeda demandada deve ser igual à moeda ofertada.

Estas seis equações seriam aceitas igualmente por adeptos da teoria quantitativa e da teoria de renda-despesa. Neste nível de abstração, não há diferença entre elas. No entanto, embora existam seis equações, há sete incógnitas: $C, I, Y, r, P, M^D, M^S$. Há uma equação faltando. Alguma destas variáveis deve ser determinada por relações fora deste sistema.\footnote{É claro que isto é falar figurativamente. Não é necessário que uma única variável seja assim determinada. O que é necessário é uma relação independente ligando algum subconjunto das sete variáveis endógenas com variáveis exógenas, e esse subconjunto poderia em princípio consistir em todas as sete variáveis.}

A diferença entre a teoria quantitativa e a teoria de renda-despesa é a condição que é adicionada para tornar as equações determinadas. A teoria quantitativa simples adiciona a equação
\begin{equation}
\frac{Y}{P} = y = y_0;
\end{equation}
isto é, a renda real é determinada fora do sistema. Com efeito, ela anexa a este sistema as equações walrasianas de equilíbrio geral, considera-as como independentes destas equações que definem os agregados, e toma como dado o valor de $Y/P$, e assim reduz este sistema a um de seis equações determinando seis incógnitas.\footnote{Esta é a essência do que tem sido chamado de dicotomia clássica. Estritamente falando, a divisão entre consumo e investimento e as taxas de câmbio entre bens atuais e futuros (o conjunto de taxas de juros "reais" ou de "propriedade") também são determinadas em um sistema walrasiano "real", que admite crescimento, razão pela qual os teóricos quantitativos tenderam a se concentrar apenas nas equações (12), (13) e (14). Deste ponto de vista, as equações (9), (10) e (11) são um resumo ou agregação ou subconjunto do sistema walrasiano.}

A teoria de renda-despesa simples adiciona a equação
\begin{equation}
P = P_0;
\end{equation}
isto é, o nível de preços é determinado fora do sistema, o que novamente reduz o sistema a um de seis equações em seis incógnitas. Ele anexa a este sistema um conjunto histórico de preços e uma estrutura institucional que é assumida como mantendo os preços rígidos ou para determinar mudanças nos preços com base no "poder de barganha" ou em algum conjunto similar de forças. Inicialmente, o conjunto de forças que determinam os preços era tratado como não sendo incorporado em qualquer corpo formal de análise econômica. Mais recentemente, os desenvolvimentos simbolizados pela "curva de Phillips" refletem tentativas de trazer a determinação de preços de volta para o corpo da análise econômica, para estabelecer um elo entre magnitudes reais e a taxa na qual os preços mudam de seu nível historicamente determinado (Phillips 1958).

Para a especialização da teoria quantitativa, dado que $Y/P = y_0$, as equações (9), (10) e (11) tornam-se um sistema autocontido de três equações em três incógnitas: $C/P, I/P,$ e $r$. Substituindo (9) e (10) em (11), temos
\begin{equation}
y_0 - f(y_0, r) = g(r),
\end{equation}
ou uma única equação que determina $r$. Seja $r_0$ este valor de $r$. Da equação (13), isso determina o valor de $M$, digamos $M_0$, que, usando a equação (14), converte (12) em
\begin{equation}
M_0 = P \cdot l(y_0, r_0),
\end{equation}
que agora determina $P$.

A equação (18) é simplesmente a equação quantitativa clássica, como pode ser visto multiplicando e dividindo o lado direito por $y_0$ e substituindo $l(y_0, r_0)/y_0$ pelo seu equivalente, $1/V$. Se retirarmos os subscritos, isso dá
\begin{equation}
M = \frac{Py}{V},
\end{equation}
ou
\begin{equation}
P = \frac{MV}{y}.
\end{equation}

Para a especialização renda-despesa, definir $P = P_0$ não dá permissão geral para uma solução sequencial. Substituindo as equações (9) e (10) na equação (11) dá
\begin{equation}
\frac{Y}{P_0} - f\left(\frac{Y}{P_0}, r\right) = g(r),
\end{equation}
uma equação em duas variáveis, $Y$ e $r$. Esta é a curva $IS$ da famosa análise $IS-LM$ de Hicks (Hicks 1937). Substituindo as equações (12) e (13) na equação (14) dá
\begin{equation}
h(r) = P_0 \cdot l\left(\frac{Y}{P_0}, r\right),
\end{equation}
uma segunda equação nas mesmas duas variáveis, $Y$ e $r$. Esta é a curva $LM$ de Hicks. A solução simultânea das duas determina $r$ e $Y$.

Alternativamente, resolva a equação (21) para $Y$ como uma função de $r$, e substitua na equação (22). Isso dá uma única equação que determina $r$ como uma função da demanda e oferta de moeda. Isso pode ser considerado como o paralelo keynesiano à equação (18), que determina $P$ como uma função da demanda e oferta de moeda.

Uma análise sequencial mais simples, fiel a muitas versões de livros didáticos da análise e ao próprio modelo simplificado de Keynes, é obtida assumindo-se que $Y/P$ não é um argumento no lado direito da equação (12) ou que a preferência absoluta pela liquidez se mantém, de modo que a equação (12) assume a forma especial:
\begin{equation}
\begin{aligned}
M^D &= 0 & \text{se } r > r_0 \\
M^D &= \infty & \text{se } r < r_0.
\end{aligned}
\tag{12a}
\end{equation}

Em qualquer um desses casos, as equações (12) ou (12a), (13) e (14) determinam a taxa de juros, $r = r_0$ (exatamente como na abordagem da teoria quantitativa as equações [9], [10] e [11] o fazem); substituindo a taxa de juros na equação (10) determina o investimento, digamos, em $I = I_0$, e na equação (9) torna o consumo uma função apenas da renda, de modo que a renda real deve então ser determinada pelo requisito de igualar a poupança ao investimento.

Se aproximarmos a função $f(Y/P, r_0)$ por uma forma linear, digamos,
\begin{equation}
\frac{C}{P} = C_0 + C_1 \frac{Y}{P},
\end{equation}
substituirmos a equação (23) na equação (11) e resolvermos para $Y/P$, obtemos
\begin{equation}
\frac{Y}{P} = \frac{C_0 + I_0}{1 - C_1},
\end{equation}
ou o simples multiplicador keynesiano, com $C_0 + I_0$ igualando o gasto autônomo e $1/(1 - C_1)$ igualando o multiplicador.

\section{A Equação Faltante: Problemas Não Resolvidos}

Essas duas abordagens têm vários elementos comuns, além de compartilharem o mesmo modelo de seis equações, que merecem destaque porque indicam quais são os principais problemas não resolvidos.

1. Ambas analisam ajustes de curto prazo em termos de deslocamentos de uma posição de equilíbrio estático para outra.
2. Ambas implicitamente consideram cada posição de equilíbrio como caracterizada por um \textit{nível} estável de preços ou produto. Nenhuma delas introduz explicitamente mudanças nos preços ou mudanças no produto na análise teórica formal. A proliferação recente de modelos de crescimento formal e a introdução ainda mais recente de mudanças monetárias neles são tentativas de preencher esta lacuna.\footnote{Alguns dos itens mais importantes são Solow (1956), Mundell (1963), Tobin (1965), Johnson (1967a, 1967b), Uzawa (1966), Sidrauski (1967a, 1967b), Levhari e Patinkin (1968) e Friedman (1969, cap. 1).}
3. Ambas consideram as taxas de juros como se ajustando instantaneamente a um novo nível de equilíbrio — na teoria quantitativa, para igualar poupança e investimento; na teoria de renda-despesa, para igualar a quantidade de moeda demandada e ofertada. Este é um retrocesso em relação ao trabalho anterior de Irving Fisher.
4. Nenhum dos modelos dá qualquer papel explícito às antecipações sobre magnitudes econômicas. A teoria de renda-despesa aproxima-se mais de fazê-lo em termos do papel que Keynes atribuiu às expectativas sobre as taxas de juros de longo prazo, que poderiam ser incorporadas na equação (12), como fizemos na equação (8). Aqui também, tem havido muito trabalho recente direcionado para preencher esta lacuna.\footnote{Alguns dos itens mais importantes são Koyck (1954), Cagan (1956), Friedman (1957), Nerlove (1958), Muth (1960), Solow (1960), Allais (1966).}
5. Ambas preenchem a equação faltante por uma suposição que não faz parte da análise teórica básica. Isso é menos flagrante, em certo sentido, para a teoria quantitativa, uma vez que pelo menos existe uma teoria econômica bem desenvolvida, resumida nas equações walrasianas de equilíbrio geral, que explica o que determina o nível de produto, de modo que as equações escolhidas para a análise podem ser vistas como um subconjunto de um sistema completo. É por isso que, à medida que se chegou a um acordo sobre a falácia da proposição (1) de Keynes, essencialmente todos os teóricos econômicos, qualquer que seja o modelo que prefiram para a análise de curto prazo, aceitam o modelo da teoria quantitativa, completado pelas equações walrasianas, como válido para o equilíbrio de longo prazo.\footnote{Veja, por exemplo, o modelo em Bailey (1962, pp. 33–36, 40–42).} A suposição de preço rígido de Keynes é, neste sentido, muito mais arbitrária. É inteiramente um \textit{deus ex machina} sem fundamento na teoria econômica. Além disso, dado que o nível de preços no longo prazo é determinado pelo modelo da teoria quantitativa, não há nenhum elo teórico entre o modelo de curto prazo e o modelo de longo prazo, nenhuma maneira de conectar um ao outro.
6. Um aspecto do ponto anterior é tão importante que merece ser afirmado explícita e separadamente. Nenhum modelo teórico tem algo a dizer sobre os fatores que determinam as proporções em que uma mudança na renda nominal irá, no curto prazo, ser dividida entre mudança de preço e mudança de produto. Uma teoria \textit{afirma} que a mudança na renda nominal será toda absorvida por mudança de preço; a outra, que será toda absorvida por mudança de produto. Na minha opinião, este é o defeito central comum das duas abordagens como teorias de mudança de curto prazo.

O resto deste artigo esboça uma possível abordagem teórica para o processo de ajuste destinada a sugerir como algumas das lacunas acima mencionadas poderiam ser preenchidas e analisa certos aspectos desta abordagem teórica em mais detalhe.\footnote{Outras partes da estrutura teórica desenvolvida mais plenamente no curso da análise empírica de algumas das questões levantadas nos outros capítulos do livro de que este artigo é abstraído.}

\section{O Processo de Ajuste}

A necessidade chave é de uma teoria para explicar (a) a divisão de curto prazo de uma mudança na renda nominal entre preços e produto, (b) o ajuste da renda nominal a uma mudança em variáveis autônomas, e (c) a transição entre esta situação de curto prazo e um equilíbrio de longo prazo descrito essencialmente pelo modelo da teoria quantitativa.

A ideia central que usarei ao esboçar a direção na qual tal teoria poderia ser desenvolvida é a distinção entre magnitudes reais e antecipadas ou, para usar uma terminologia que não precisa ser idêntica, mas que tratarei para este propósito como se fosse, entre magnitudes medidas e permanentes. Numa posição de equilíbrio de longo prazo, todas as antecipações são realizadas, de modo que as magnitudes reais e antecipadas, ou medidas e permanentes, são iguais.\footnote{Note que a igualdade de magnitudes reais e antecipadas é uma condição necessária, mas não suficiente, para uma posição de equilíbrio de longo prazo. Em princípio, magnitudes reais e antecipadas poderiam ser iguais ao longo de um caminho de ajuste entre uma posição de equilíbrio e outra. A proposição correspondente é mais complicada para magnitudes medidas e permanentes e depende da definição precisa desses termos. No entanto, uma vez que consideraremos um caso especial no qual a condição estabelecida é tratada como necessária e suficiente para o equilíbrio de longo prazo, estas complicações podem ser contornadas.}

Considerarei o equilíbrio de longo prazo como determinado pelo modelo da teoria quantitativa anterior mais as equações walrasianas de equilíbrio geral. Num enunciado completo, o modelo anterior deveria ser expandido incluindo a riqueza nas funções de consumo e preferência pela liquidez, e o estoque de capital na função de investimento, e permitindo o crescimento constante do produto e dos preços.

Considerarei o equilíbrio de curto prazo como determinado por um processo de ajuste no qual a taxa de ajuste de uma variável é uma função da discrepância entre os valores medidos e os antecipados dessa variável ou de sua taxa de variação, bem como, talvez, de outras variáveis ou de suas taxas de variação. Finalmente, deixarei pelo menos algumas variáveis antecipadas serem determinadas por um processo de retroalimentação a partir de valores observados passados.

\subsection{Divisão de uma Mudança na Renda Nominal entre Preços e Produto}

Parece plausível que a divisão de uma mudança na renda nominal entre preços e produto dependa de dois fatores principais: antecipações sobre o comportamento dos preços — este é o fator de inércia enfatizado por Keynes — e o nível atual de produto ou emprego comparado com o nível de pleno emprego (permanente) de produto ou emprego — este é o fator de demanda-oferta enfatizado pelos teóricos quantitativos. Podemos expressar isto na forma geral como:
\begin{equation}
\frac{dP}{dt} = f\left[ \frac{dY}{dt}, \left(\frac{dP}{dt}\right)^*, \left(\frac{dy}{dt}\right)^*, y, y^* \right],
\end{equation}
\begin{equation}
\frac{dy}{dt} = g\left[ \frac{dY}{dt}, \left(\frac{dP}{dt}\right)^*, \left(\frac{dy}{dt}\right)^*, y, y^* \right],
\end{equation}
onde um asterisco anexado a uma variável denota o valor antecipado dessa variável e onde a forma das equações (25) e (26) deve ser consistente com a identidade
\begin{equation}
Y = Py,
\end{equation}
de modo que apenas uma das equações (25) e (26) seja independente.

Para ilustrar, uma versão linearizada específica das equações (25) e (26) poderia ser
\begin{equation}
\frac{d \log P}{dt} = \left( \frac{d \log P}{dt} \right)^* + \alpha \left[ \frac{d \log Y}{dt} - \left( \frac{d \log Y}{dt} \right)^* \right] + \gamma [\log y - (\log y)^*];
\end{equation}
\begin{equation}
\frac{d \log y}{dt} = \left( \frac{d \log y}{dt} \right)^* + (1 - \alpha) \left[ \frac{d \log Y}{dt} - \left( \frac{d \log Y}{dt} \right)^* \right] - \gamma [\log y - (\log y)^*].
\end{equation}

A soma de ambas é exatamente o logaritmo da equação (27), diferenciado em relação ao tempo, desde que as variáveis antecipadas também satisfaçam uma identidade correspondente,\footnote{Isso também explica por que $[(d \log y)/dt]^*$ não parece aparecer explicitamente na equação (28), ou $[(d \log P)/dt]^*$ na equação (29), visto que elas entram nas equações (25) e (26). Elas estão implicitamente incluídas em $[(d \log Y)/dt]^*$.} de modo que as equações satisfaçam as condições especificadas.

A suposição simples da teoria quantitativa, de que toda a mudança na renda está nos preços, e que o produto está sempre no seu nível permanente, é obtida fixando $\alpha = 1$ e $\gamma = \infty$. Um valor infinito de $\gamma$ corresponde a "preços perfeitamente flexíveis" e assegura que $y = y^*$. O valor unitário de $\alpha$ assegura que os preços absorvam qualquer mudança na renda nominal, de modo que a renda real cresça na sua taxa de crescimento de longo prazo.\footnote{Com $\gamma$ infinito, e $\log y = \log y^*$, a expressão final nas equações (28) e (29) é $\infty \cdot 0$, ou tecnicamente indeterminada. O produto pode ser tomado como sendo zero em geral, exceto possivelmente para alguns pontos isolados nos quais $\log y$ se desvia de $\log y^*$, um desvio fechado instantaneamente por taxas infinitas de mudança em $\log P$ e $\log y$.}

A suposição keynesiana simples, de que toda a mudança na renda está no produto, enquanto houver desemprego, e toda nos preços, uma vez que haja pleno emprego, é obtida fixando $[(d \log P)/dt]^* = 0$, e $\alpha = \gamma = 0$ para $y < y^*$, e então mudando para a especificação da teoria quantitativa $\alpha = 1, \gamma = \infty$ para $y \ge y^*$. O valor zero de $[(d \log P)/dt]^*$ assegura que as antecipações sejam de preços estáveis e, combinado com os valores zero de $\alpha$ e $\gamma$, que $(d \log P)/(dt) = 0$. Seria um pouco mais geral, e talvez mais consistente com o espírito, em vez da letra da análise de Keynes, e ainda mais com a de seus seguidores modernos, deixar $[(d \log P)/(dt)]^*$ diferir de zero enquanto mantém $\alpha = \gamma = 0$ para $y < y^*$. Isso introduziria o tipo de rigidez de preço relevante para a análise de curto período de Keynes, mas ainda poderia ser considerado como capturando o fenômeno que seus seguidores modernos enfatizaram como inflação de custo.

Nas suas formas gerais, as equações (28) e (29) não especificam por si mesmas o caminho dos preços ou do produto a partir de qualquer posição inicial. Além disso, precisamos saber como os valores antecipados são formados. Presumivelmente, estes são afetados pelo curso dos eventos, de modo que, em resposta a um distúrbio que produz uma discrepância entre os valores reais e antecipados das variáveis, há um efeito de retroalimentação que aproxima novamente os valores reais e antecipados (ver abaixo). Se este processo prossegue rapidamente, então os ajustes transitórios definidos pelas equações (28) e (29) são de pouca significância. A análise relevante é a análise que conecta as variáveis com asterisco.

\subsection{Ajuste de Curto Prazo da Renda Nominal}

Para a teoria monetária, a questão chave é o processo de ajuste a uma discrepância entre a quantidade nominal de moeda demandada e a quantidade nominal de moeda ofertada. Tal discrepância poderia surgir de uma mudança na oferta de moeda (um deslocamento na função de oferta) ou de uma mudança na demanda de moeda (um deslocamento na função de demanda). O \textit{insight} chave da abordagem da teoria quantitativa é que tal discrepância se manifestará principalmente em gastos tentados, e daí na taxa de mudança da renda nominal. Dito de outra forma, os detentores de moeda não podem determinar a quantidade nominal de moeda (embora suas reações possam introduzir efeitos de retroalimentação que afetarão a quantidade nominal de moeda), mas eles podem tornar a velocidade tudo o que desejarem.

O que, nesta visão, fará com que a taxa de mudança na renda nominal se desvie de seu valor permanente? Qualquer coisa que produza uma discrepância entre a quantidade nominal de moeda demandada e a quantidade ofertada, ou entre as duas taxas de variação da moeda demandada e da moeda ofertada. Em forma geral
\begin{equation}
\frac{dY}{dt} = f\left[ \left(\frac{dY}{dt}\right)^*, \frac{dM^S}{dt}, \frac{dM^D}{dt}, M^S, M^D \right],
\end{equation}
onde $M^S$ refere-se à moeda ofertada, $M^D$ à moeda demandada, e os dois símbolos são usados para indicar que as duas não são necessariamente iguais. Ou seja, a equação (30) substitui a equação de ajuste (14), $M^D = M^S$, do modelo simples.

Para ilustrar, uma versão linearizada particular da equação (30) seria
\begin{equation}
\frac{d \log Y}{dt} = \left( \frac{d \log Y}{dt} \right)^* + \Psi \left( \frac{d \log M^S}{dt} - \frac{d \log M^D}{dt} \right) + \Phi (\log M^S - \log M^D).
\end{equation}

Diferente das equações (28) e (29), os dois termos finais de ajuste no lado direito não incluem explicitamente quaisquer magnitudes com asterisco. Mas implicitamente o fazem. O montante de moeda demandada dependerá da renda permanente antecipada e dos preços, bem como da taxa antecipada de variação dos preços.

O caso especial do modelo simples da teoria quantitativa pode ser obtido a partir da equação (31) fixando $(d \log M^D)/(dt) = [(d \log Y)/(dt)]^*, \Psi = 1$ e $\Phi = \infty$. A primeira condição é simplesmente que a velocidade de longo prazo é constante, a terceira ($\Phi = \infty$) assegura que $M^S = M^D$ (equação [14]), e a segunda ($\Psi = 1$), que $(d \log Y)/(dt) = (d \log M^S)/(dt)$; isto é, uma mudança na oferta de moeda é refletida imediata e proporcionalmente na renda nominal.

Especificar o modelo keynesiano simples em termos da equação (31) é mais difícil porque o ajuste inicial naquela abordagem é assumido como sendo através das taxas de juros. Uma contraparte da equação (31) seria a taxa de variação nas taxas de juros no lado esquerdo. Uma mudança na taxa de juros, por sua vez, é assumida para influenciar a renda através do seu efeito no investimento. Outra contraparte da equação (31) teria, no lado direito, discrepâncias entre as taxas reais e antecipadas de juros e nível de gastos autônomos.

Em seus termos gerais, a equação (31) permite mudanças tanto na oferta quanto na demanda por moeda. Ela também permite implicitamente as forças enfatizadas por Keynes, deslocamentos no investimento ou outros gastos autônomos, através do efeito de tais mudanças em $M^S$ e $M^D$. Por exemplo, um aumento autônomo na demanda de investimento tenderá a elevar as taxas de juros. O aumento nas taxas de juros tenderá a reduzir $M^D$, introduzindo uma discrepância em um ou ambos os termos entre parênteses no lado direito da equação (31), o que fará com que $(d \log Y)/(dt)$ exceda $[(d \log Y)/(dt)]^*$.

\subsection{Funções de Demanda e Oferta de Moeda}

Como este comentário indica, para completar a teoria do processo de ajuste, é necessário especificar as funções que conectam $M^D$ e $M^S$ com outras variáveis no sistema, e também fornecer relações que determinem quaisquer variáveis adicionais — como taxas de juros — que entrem nestas funções. As seções 3 e 4 acima discutem a demanda e a oferta de moeda que considero relevantes para este propósito, de modo que apenas alguns breves comentários suplementares são necessários para os propósitos presentes.

Primeiro, em muito do meu trabalho empírico, tomei $M^S$ como uma variável autônoma e não a incorporei na análise de qualquer retroalimentação de outros ajustes. Uma razão importante pela qual o fiz é meu julgamento de que a função de oferta variou grandemente de tempos em tempos.

Segundo, na notação que utilizei nesta seção, as variáveis $y$ e $(1/P)(dP/dt)$ na equação (7) deveriam ter asteriscos anexados a elas.

Terceiro, a função que especifica $M^D$ poderia, em princípio, incluir um componente transitório. Ou seja, não há nada inconsistente com a teoria aqui esboçada em distinguir entre uma demanda por moeda de curto prazo e uma de longo prazo, como alguns autores fizeram (Heller 1965; Chow 1966; Konig 1968).

\subsection{Determinação das Taxas de Juros}

Dado que as taxas de juros entram na função de demanda por moeda (equação [7]) e também, presumivelmente, na função de oferta, um modelo completo deve especificar os fatores que as determinam. Nosso modelo de longo prazo determina seus valores permanentes. O que é necessário é uma análise do processo de ajuste para as taxas de juros comparável àquela para preços e renda nominal discutida acima — desde que, como parece razoável, as taxas de juros medidas, bem como as taxas de juros permanentes, entrem nas funções de demanda e oferta de moeda.

Eu mesmo não desenvolvi a teoria formal deste processo de ajuste em qualquer detalhe. O único aspecto que considerei é o efeito de mudanças em $M^S$ sobre as taxas de juros.\footnote{Ver capítulo 7 no livro futuro do qual este artigo foi abstraído.} Nessa análise, tratei as taxas de juros como se ajustando muito rapidamente para limpar o mercado para fundos de empréstimos, sendo a oferta de fundos de empréstimos possivelmente ligada a mudanças em $M^S$, e a demanda e a oferta de fundos de empréstimos expressas como uma função da taxa de juros nominal conforme dependente de $Y$ e $[(1/P)(dP/dt)]^*$ junto com outras variáveis.

Em parte do meu trabalho empírico, tratei as taxas de juros como exógenas.

\subsection{Determinação de Valores Antecipados}

A transição entre o processo de ajuste de curto prazo e o equilíbrio de longo prazo é produzida por um ajuste de valores antecipados aos valores medidos de tal forma que, para um sistema estável, um único distúrbio estabelecerá discrepâncias que serão eliminadas com o tempo. Para colocar isto em termos gerais, devemos ter:
\begin{equation}
\left[ \frac{d \log P}{dt} (t) \right]^* = f\left[ \frac{d \log P}{dt} (T) \right],
\end{equation}
\begin{equation}
\left[ \frac{d \log Y}{dt} (t) \right]^* = g\left[ \frac{d \log Y}{dt} (T) \right],
\end{equation}
\begin{equation}
y^*(t) = h[y(T)],
\end{equation}
\begin{equation}
P^*(t) = j[P(T)],
\end{equation}
onde $t$ representa um ponto particular no tempo e $T$ um vetor de todas as datas anteriores a $t$.

Um distúrbio do equilíbrio de longo prazo, digamos, introduz discrepâncias nos dois termos entre parênteses no lado direito da equação (31). Isso fará com que a taxa de mudança na renda nominal se desvie de seu valor permanente, o que, através das equações (28) e (29), produz desvios na taxa de mudança de preço e na mudança de produto de seus valores permanentes. Estes podem, por sua vez, reentrar na equação (31) através de mudanças em $k$ se o fizerem, eles irão, através das equações (32)–(35), produzir mudanças nos valores antecipados que irão, mais cedo ou mais tarde e talvez após um processo de reação cíclica, eliminar as discrepâncias entre os valores medidos e os permanentes.

Estas equações de antecipação são num sentido muito gerais, noutro, muito especiais. Elas exigem que as antecipações sejam determinadas inteiramente pela história passada da variável particular em questão, não por outra história passada ou outras alternativas observadas atualmente. Estas equações negam qualquer papel "autônomo" às antecipações. Estas equações, ou alternativas preferíveis a elas, não estão diretamente relacionadas às questões monetárias que são a principal preocupação deste artigo, razão pela qual as tratei tão sumariamente. Sua única função aqui é fechar o sistema.

Um problema sutil neste tipo de estrutura, no qual identificamos a ausência de uma discrepância entre os valores reais e os antecipados como definindo o equilíbrio de longo período, é assegurar que as relações de retroalimentação definidas pelas (32)–(35), bem como as outras funções, sejam consistentes com o sistema expandido de equações walrasianas que especificam os valores de equilíbrio de longo prazo. Pelo menos alguns valores são implicitamente determinados de duas formas: por uma relação de retroalimentação tal como as equações (32) e (35), e pelo sistema de equações de equilíbrio de longo prazo. O problema é assegurar que no equilíbrio de longo prazo estas duas determinações não conflitem.

No meu trabalho empírico, geralmente usei uma forma particular de função de antecipação, a saber, aquela que define os valores antecipados como uma média ponderada decrescente de valores observados passados. Por exemplo, uma forma específica da equação (34) é:
\begin{equation}
y^*(t) = \beta \int_{T=-\infty}^t e^{(\beta - \alpha)(t - T)} y(T) dT,
\end{equation}
onde $\alpha$ e $\beta$ são parâmetros, $\alpha$ definindo a taxa de crescimento de longo prazo, e $\beta$ a velocidade de ajuste das antecipações às experiências (Friedman 1957, pp. 142–47).

\section{Uma Ilustração}

Pode ajudar a esclarecer a natureza geral desta abordagem teórica se a aplicarmos a um distúrbio monetário hipotético.

Comecemos com uma situação de equilíbrio pleno com preços estáveis e pleno emprego e com o produto crescendo a, digamos, 3 por cento ao ano. Por simplicidade, assuma que a elasticidade-renda da demanda de moeda é unitária, de modo que a quantidade de moeda também cresça à taxa de 3 por cento ao ano. Assuma também que a moeda é moeda fiduciária integralmente não remunerada e que sua quantidade pode ser tomada como autônoma.

Suponha que haja um deslocamento no tempo $t = t_0$ na taxa de crescimento da quantidade de moeda de 3 por cento ao ano para, digamos, 8 por cento ao ano e que esta nova taxa de crescimento seja mantida indefinidamente. A Figura 1 mostra a trajetória temporal do estoque monetário antes e depois do tempo $t_0$. Estas figuras não são desenhadas rigorosamente à escala. Por ênfase, elas exageram a diferença nas inclinações das linhas antes e depois de $t_0$.

\subsection{Equilíbrio de Longo Prazo}

Perguntemo-nos primeiro qual será a solução de equilíbrio de longo prazo. Claramente, após o ajuste total, a renda nominal estará crescendo a 8 por cento ao ano. Se, por enquanto, negligenciarmos qualquer efeito desta mudança monetária sobre o produto real e a taxa de crescimento do produto, isso significa que os preços estariam subindo a 5 por cento ao ano. Pode parecer, portanto, que se a trajetória de equilíbrio da renda nominal fosse simplesmente multiplicada pela trajetória da quantidade de moeda na figura 1 (redesenhada como a linha sólida mais a tracejada na fig. 2). Mas não é o caso. Com os preços subindo à taxa de 5 por cento ao ano e, em equilíbrio, com esse aumento de preço totalmente antecipado por todos, agora é mais dispendioso manter moeda. Como resultado, a equação (7) indicaria um declínio na quantidade real de moeda demandada em relação à renda, isto é, um aumento na velocidade desejada. Este aumento seria alcançado por um aumento na renda nominal acima do montante necessário para corresponder ao aumento na quantidade nominal de moeda. A trajetória de equilíbrio da renda nominal seria como a linha sólida na figura 2 em vez da linha tracejada.

\begin{figure}[h]
\centering
% Placeholder para Fig 1
\textbf{[Figura 1: Trajetória temporal do estoque monetário antes e depois do tempo $t_0$]}
\end{figure}

Se o produto real de equilíbrio e a taxa de crescimento do produto real fossem inalterados pela mudança monetária, como assumi até agora, a trajetória de equilíbrio dos preços seria a mesma da renda nominal, exceto que teria uma inclinação de 3 por cento ao ano a menos, para permitir o crescimento da renda real. No entanto, o nível de equilíbrio do produto real não seria inafetado pela mudança monetária. O efeito exato depende de quão exatamente o produto real é medido, em particular se ele inclui ou exclui os serviços não pecuniários da moeda. Se os incluir, como em princípio deveria, então o nível de produto real será menor após a mudança monetária do que antes. Ele será menor por duas razões: primeiro, o custo mais alto de manter saldos de caixa levará os produtores a substituir o caixa por outros recursos, o que diminuirá a eficiência produtiva; segundo, o fluxo de serviços não pecuniários da moeda será reduzido (Friedman 1969, pp. 14–15). Por ambas as razões, o nível de preço do produto terá de subir mais do que a renda nominal — uma linha sólida e uma linha tracejada como aquelas para a renda nominal na figura 2 estariam mais distantes verticalmente para os preços do produto final do que para a renda nominal.

\begin{figure}[h]
\centering
% Placeholder para Fig 2
\textbf{[Figura 2: Trajetória de equilíbrio da renda nominal antes e depois do tempo $t_0$]}
\end{figure}

É mais difícil ser preciso sobre a taxa de crescimento de equilíbrio, visto que ela depende do modelo de crescimento particular. O que está claro é que o estoque agregado de capital físico (não monetário), incluindo a moeda, será menor em relação ao capital humano, mas que o estoque agregado de capital físico (não monetário) será maior, de modo que o rendimento real do capital (essencialmente nosso $r_e$ da equação [7]) será menor. A taxa de juros nominal (o $r_b$ da equação [7]) será igual a este rendimento real mais a taxa de variação dos preços, logo será maior. Se estas mudanças têm algum efeito na taxa de crescimento do produto real, elas tenderão a reduzi-lo, de modo que o nível de equilíbrio do preço do produto final não apenas será maior em relação ao seu valor inicial do que o nível de equilíbrio da renda nominal; ele também poderá subir mais rapidamente (Stein 1966; Johnson 1967a, 1967b; Marty 1968). Por simplicidade, negligenciarei esta possibilidade e assumirei que a taxa de crescimento de equilíbrio dos preços é de 5 por cento ao ano.

\subsection{O Processo de Ajuste}

Até aqui para a posição de equilíbrio. O que dizer sobre o processo de ajuste?
Esta descrição da posição de equilíbrio já nos diz uma coisa sobre o processo de ajuste. Para produzir o deslocamento na trajetória de equilíbrio da renda nominal da tracejada para a sólida, a renda nominal e os preços devem subir durante algum período a uma taxa mais rápida do que a taxa final de equilíbrio — a uma taxa mais rápida do que 8 por cento ao ano para a renda nominal e 5 por cento ao ano para os preços. Deve haver, portanto, uma reação cíclica, um "overshooting", na taxa de mudança da renda nominal e dos preços, embora não necessariamente nos seus níveis.

Como este processo de ajuste será refletido no meu esboço teórico do processo de ajuste? O deslocamento em $(d \log M^S)/(dt)$ no tempo $t_0$ de 3 por cento para 8 por cento introduz uma discrepância de sinal positivo no segundo termo entre parênteses da equação (31), enquanto deixa inicialmente o terceiro termo entre parênteses inalterado. Como resultado, $(d \log Y)/(dt)$ aumentará, excedendo $[(d \log Y)/(dt)]^*$, que, visto neste processo de transição como um valor antecipado em vez de como um valor de equilíbrio de longo prazo, é inalterado em relação ao valor de equilíbrio anterior de longo prazo. Quão rapidamente a taxa de crescimento da renda nominal sobe depende em parte do valor de $\Psi$, o coeficiente que indica a velocidade de ajuste, e em parte da função de demanda por moeda. Se este último depende apenas de valores antecipados (isto é, se todas as variáveis na equação [7] têm asteriscos), $(d \log M^D)/(dt)$ será inicialmente inalterado, então tudo dependerá de $\Psi$, que pode ter qualquer valor, de zero, significando nenhum ajuste, até um valor maior que a unidade, significando que a renda nominal subiria inicialmente por mais de 5 por cento ao ano.\footnote{O modelo esboçado brevemente nos dois parágrafos finais de Friedman (1959) implicitamente tem um valor inicial de $\Psi$ que é muito maior que a unidade.}

Qualquer que seja a taxa de aumento da renda nominal, ela será dividida em um aumento nos preços e no produto, de acordo com as equações (28) e (29). Se $\alpha$ for menor que a unidade, tanto o produto real quanto os preços começarão a subir, suas taxas relativas dependendo do tamanho de $\alpha$.

O aumento dos preços e da renda nominal começará a afetar as taxas antecipadas de variação, através das equações (32) e (33), realimentando (31) e (28) e (29).

Tudo isto acontece no tempo $t_0$, sem efeito sobre os níveis de quaisquer das variáveis. No entanto, à medida que o processo continua, os níveis começam a ser afetados. Na equação (31), $\log M^S$ passa a exceder $\log M^D$, de modo que o segundo termo da equação (31) adiciona à pressão ascendente sobre $(d \log Y)/(dt)$, fazendo com que a expansão da renda nominal se acelere. Nas equações (28) e (29), $\log y$ passa a exceder $\log y^*$, aumentando assim a fração da renda nominal absorvida pelos preços e reduzindo a fração absorvida pelo produto. Os níveis alterados de $y$ e $P$ realimentam as equações (34) e (35) e assim começam a alterar $y^*$ e $P^*$.

As mudanças em todas as variáveis agora começam a afetar as funções de demanda por moeda, tanto diretamente, à medida que estas variáveis entram nas funções de demanda, quanto indiretamente, à medida que afetam outras variáveis, como taxas de juros, que por sua vez entram nas funções de demanda. À medida que $(d \log M^D)/(dt)$ e $M^D$ na equação (31) começam a mudar. O processo, é claro, será finalmente completado quando as variáveis medidas relevantes forem todas iguais às suas contrapartes permanentes e estas forem iguais aos valores de equilíbrio de longo prazo discutidos acima.

É impossível levar muito mais longe esta afirmação verbal da solução de um sistema incompletamente especificado de equações diferenciais simultâneas. O caminho exato de ajuste depende de como os elementos faltantes do sistema são especificados e dos valores numéricos dos parâmetros, mas talvez isto seja suficiente para dar o sabor do tipo de processo de ajuste que eles geram, e para indicar por que este processo é necessariamente cíclico.

Qual é o reflexo nestas equações do ponto feito no segundo parágrafo desta seção, a saber, que $(d \log Y)/(dt)$ e $(d \log P)/(dt)$ devem, durante a transição, ser em média maiores do que os seus valores finais de equilíbrio de longo prazo? Considere a equação (31). Suponha que durante um período o valor médio de $(d \log Y)/(dt)$ e $(d \log P)/(dt)$ tenha sido 8 por cento ao ano e 5 por cento ao ano, respectivamente. Suponha que as funções de antecipação (32) e (35) tenham refletido integralmente estes valores antecipados. Então, como vimos, embora $M^S$ tivesse crescido à taxa de 5 por cento ao ano, $M^D$ não o teria; de modo que o termo final na equação (31) não seria zero, mesmo que o termo do meio no lado direito pudesse ser. Logo, $(d \log Y)/(dt)$ excederia $[(d \log Y)/(dt)]^*$, que por suposição é a sua taxa de crescimento de longo prazo; portanto, o equilíbrio pleno não teria sido atingido.

A Figura 3 resume vários caminhos de ajuste possíveis de $(d \log Y)/(dt)$ consistentes com a teoria esboçada. A característica comum de todos eles é que a área acima da linha de 8 por cento deve exceder a área abaixo. Em princípio, é claro, outros caminhos ainda são possíveis. Por exemplo, é conceitualmente possível para o ajuste ser explosivo em vez de amortecido. Restringir-nos a caminhos amortecidos é um julgamento empírico.

\begin{figure}[h]
\centering
% Placeholder para Fig 3
\textbf{[Figura 3: Possíveis caminhos de ajuste da taxa de variação na renda nominal]}
\end{figure}

\section{Conclusão}

Ao concluir esta discussão de uma estrutura teórica, pode valer a pena afirmar que não se trata de uma estrutura especial para mim ou para aqueles economistas que veem a operação da economia nos termos da teoria quantitativa. Sem dúvida, outros economistas expandiriam a estrutura de forma diferente, enfatizariam diferentes partes dela, elaborariam pontos que eu passei por cima, e passariam por cima de pontos que eu elaborei. Mas quase todos os economistas aceitariam a estrutura, e isso é verdade mesmo, acredito, da parte menos exaustiva, o esboço do processo de ajuste na seção precedente.

Um propósito de estabelecer esta estrutura, como fiz, é documentar minha crença de que as diferenças básicas entre os economistas são empíricas, não teóricas: quão importantes são as mudanças na oferta de moeda comparadas com mudanças na demanda por moeda? São transações variáveis ou variáveis de ativos as mais importantes na determinação da demanda por moeda? Quão elástica é a demanda por moeda em relação às taxas de juros? Com relação à taxa de variação nos preços? Quando ocorrem mudanças na demanda ou na oferta, produzem discrepâncias entre a quantidade de moeda que o público detém e a quantidade que deseja deter, quão rapidamente essas discrepâncias tendem a ser eliminadas? O ajuste atinge principalmente os preços ou principalmente as quantidades? O processo de ajuste é cíclico ou assintótico? O ajuste a mudanças abruptas em curtos períodos é diferente em espécie ou apenas em grau do ajuste a mudanças mais lentas em períodos mais longos? Quanto tempo leva para as pessoas alterarem suas antecipações à luz da experiência?

Muita da controvérsia que girou em torno do papel da moeda nos assuntos econômicos reflete, na minha opinião, diferentes respostas implícitas ou explícitas a estas questões empíricas. A razão pela qual tais diferenças foram capazes de persistir é, acredito, que o ajuste total a distúrbios monetários leva um tempo muito longo e afeta muitas magnitudes econômicas. Se o ajuste fosse rápido, imediato e mecânico, como alguns teóricos quantitativos anteriores podem ter acreditado, ou, mais provavelmente, como foi atribuído a eles por seus críticos, o papel da moeda seria clara e nitidamente gravado nas figuras imperfeitas que têm estado disponíveis. Mas, se o ajuste for lento, retardado e sofisticado, então a evidência bruta pode ser enganosa, e uma análise mais sutil do registro pode ser necessária para desembaraçar o que é sistemático do que é aleatório e errático. Esse, e não a elaboração da teoria, tem sido o objetivo primário dos estudos meus e de meus associados.

\section*{Referências}

Allais, M. "A Restatement of the Quantity Theory of Money." \textit{A.E.R.} (Dezembro 1966), pp. 1123–57.

Bailey, Martin J. \textit{National Income and the Price Level}. Nova York: McGraw-Hill, 1962.

Baumol, W. J. "The Transactions Demand for Cash: An Inventory Theoretic Approach." \textit{Q.J.E.} (Novembro 1952), pp. 545–56.

Brainard, William. "Financial Institutions and a Theory of Monetary Control." In \textit{Financial Markets and Economic Activity}. Monografia da Fundação Cowles 21. Nova York: Wiley, 1967.

Brunner, Karl, e Meltzer, Allan H. "Predicting Velocity." \textit{J. Finance} (Maio 1963), pp. 319–34.

Cagan, Phillip. "The Monetary Dynamics of Hyperinflation." In \textit{Studies in the Quantity Theory of Money}, editado por Milton Friedman. Chicago: Univ. Chicago Press, 1956.

———. \textit{Determinants and Effects of Changes in the Stock of Money, 1875–1960}. Nova York: Columbia Univ. Press (para Nat. Bur. Econ. Res.), 1965.

Chow, Gregory C. "On the Long-Run and Short-Run Demand for Money." \textit{J.P.E.} 74 (Abril 1966): 111–31.

Culbertson, J. M. "United States Monetary History: Its Implications for Monetary Theory." \textit{Nat. Banking Rev.} (Março 1964), pp. 372–75.

Fisher, Irving. \textit{The Purchasing Power of Money}. Nova York: Macmillan, 1911.

———. "Money, Prices, Credit, and Banking." \textit{A.E.R.} (Junho 1919), pp. 407–9.

Friedman, Milton. "The Quantity Theory of Money—a Restatement." In \textit{Studies in the Quantity Theory of Money}, editado por M. Friedman. Chicago: Univ. Chicago Press, 1956. Reimpresso em Friedman (1969).

———. \textit{A Theory of the Consumption Function}. Princeton, N.J.: Princeton Univ. Press (para Nat. Bur. Econ. Res.), 1957.

———. "The Supply of Money and Changes in Prices and Output." In \textit{The Relationship of Prices to Economic Stability and Growth}, pp. 241–56. U.S. Congress, Joint Economic Committee, Compendium, 1958. Reimpresso em Friedman (1969).

———. "The Demand for Money: Some Theoretical and Empirical Results." \textit{J.P.E.} 67 (Agosto 1959): 327–51. Reimpresso como Occasional Paper 68 (New York: Nat. Bur. Econ. Res., 1959) e em Friedman (1969).

———. \textit{Price Theory}. Chicago: Aldine, 1962.

Friedman, Milton. "Interest Rates and the Demand for Money." \textit{J. Law and Econ.} (Outubro 1966), pp. 71–85. Reimpresso em Friedman (1969).

———. "The Monetary Theory and Policy of Henry Simons." \textit{J. Law and Econ.} (Outubro 1967), pp. 1–13. Reimpresso em Friedman (1969).

———. "Money, Quantity Theory." In \textit{International Encyclopedia of the Social Sciences}, pp. 432–47. Nova York: Macmillan and Free Press, 1968. (a)

———. "The Role of Monetary Policy." \textit{A.E.R.} (Março 1968), pp. 1–17. (b) Reimpresso em Friedman (1969).

———. \textit{The Optimum Quantity of Money and Other Essays}. Chicago: Aldine, 1969.

Friedman, Milton, e Schwartz, Anna. "Money and Business Cycles." \textit{Rev. Econ. and Statis.}, supp. (Fevereiro 1963), pp. 32–64. Reimpresso em Friedman (1969). (a)

———. \textit{A Monetary History of the United States, 1867–1960}. Princeton, N.J.: Princeton Univ. Press, (para Nat. Bur. Econ. Res.), 1963. (b)

———. \textit{Monetary Statistics of the United States}. Nova York: Columbia Univ. Press (para Nat. Bur. Econ. Res.), 1970.

Goldfeld, Stephen M. \textit{Commercial Bank Behavior and Economic Activity}. Amsterdã: North-Holland, 1966.

Gramley, Lyle, e Chase, S. B., Jr. "Time Deposits in Monetary Analysis." \textit{Federal Reserve Bull.} 51 (Outubro 1965): 1380–1406.

Haberler, Gottfried. \textit{Prosperity and Depression}. 3ª ed. Genebra: Sociedade das Nações, 1941.

Heller, H. R. "The Demand for Money: The Evidence from the Short-Run Data." \textit{Q.J.E.} (Maio 1965), pp. 291–303.

Hendershott, Patric H. \textit{The Neutralized Money Stock}. Homewood, Ill.: Irwin, 1968.

Hester, Donald, e Tobin, James, eds. \textit{Financial Markets and Economic Activity}. Monografia da Fundação Cowles 21. Nova York: Wiley, 1967.

Hicks, J. R. "Mr. Keynes and the Classics: A Suggested Interpretation." \textit{Econometrica} 5 (Abril 1937): 147–59. Reimpresso em \textit{Readings in the Theory of Income Distribution}, editado por W. Fellner e B. F. Haley. Homewood, Ill: Irwin, 1951.

Holzman, Franklyn D., e Bronfenbrenner, Martin. "Survey of Inflation Theory." \textit{A.E.R.} (Setembro 1963), pp. 593–661.

Johnson, Harry G. "The General Theory after Twenty-Five Years." \textit{A.E.R.} (Maio 1961), pp. 1–17.

———. "The Neo-classical One-Sector Growth Model: A Geometrical Exposition and Extension to a Monetary Economy." In \textit{Essays in Monetary Economics}. Londres: Allen & Unwin, 1967. (a)

———. "Neutrality of Money in Growth Models: A Reply." \textit{Economica} (Fevereiro 1967), pp. 73–74. (b)

Keynes, John M. \textit{The General Theory of Employment, Interest and Money}. Londres: Macmillan, 1936.

Konig, H. "Demand Function, Short-Run and Long-Run Function, and the Distributed Lag." \textit{Zeitschrift Gesamte Staatswissenschaft} (Fevereiro 1968), pp. 124 ss.

Koyck, L. M. \textit{Distributed Lags and Investment Analysis}. Amsterdã: North-Holland, 1954.

Laurent, Robert. "Currency Transfers by Denominations." Tese de doutorado, Univ. Chicago, 1969.

Leijonhufvud, Axel. \textit{On Keynesian Economics and the Economics of Keynes}. Londres: Oxford Univ. Press, 1968.

Levhari, D., e Patinkin, D. "The Role of Money in a Simple Growth Model." \textit{A.E.R.} (Setembro 1968), pp. 713–53.

Martin, P. W. \textit{The Flaw in the Price System}. Londres: King, 1924.

Marty, Alvin. "The Optimal Rate of Growth of Money." \textit{J.P.E.} 76, pt. 2 (Julho–Agosto 1968): 860–73.

Meltzer, Allan H. "The Demand for Money: The Evidence from the Time Series." \textit{J.P.E.} 71 (Junho 1963): 219–46.

———. "Monetary Theory and Monetary History." \textit{Schweizerische Zeitschrift Volkswirtschaft und Statistik}, no. 4 (1965), pp. 409–22.

Mitchell, W. C. \textit{Business Cycles}. Nova York: Nat. Bur. Econ. Res., 1927.

Mundell, Robert A. "A Fallacy in the Interpretation of Macro-economic Equilibrium." \textit{J.P.E.} 71 (Fevereiro 1963): 61–66.

Muth, J. F. "Optimal Properties of Exponentially Weighted Forecasts." \textit{J. American Statis. Assoc.} (Junho 1960), pp. 299–306.

Nerlove, Marc. \textit{Distributed Lags and Demand Analysis}. Agriculture Handbook no. 141. Washington: Dept. Agriculture, 1958.

Okun, Arthur M. "Comment." \textit{Rev. Econ. and Statis.}, supp. (Fevereiro 1963), pp. 72–77.

Patinkin, Don. "Price Flexibility and Full Employment." In \textit{Readings in Monetary Theory}, editado por F. A. Lutz e L. W. Mints. Homewood, Ill.: Irwin, 1951. Versão revisada de um artigo originalmente publicado em \textit{A.E.R.} (Setembro 1948), pp. 543–64.

Phillips, A. W. "The Relation between Unemployment and the Rate of Change of Money Wage Rates in the United Kingdom, 1861–1957." \textit{Economica} (Novembro 1958), pp. 283–99.

Pigou, A. C. "The Value of Money." \textit{Q.J.E.} (Novembro 1917), pp. 38–65. Reimpresso em \textit{Readings in Monetary Theory}, editado por F. A. Lutz e L. W. Mints. Homewood, Ill.: Irwin, 1951.

———. "Economic Progress in a Stable Environment." \textit{Economica}, n.s. (Agosto 1947), pp. 180–88.

Sidrauski, M. "Rational Choice and Patterns of Growth in a Monetary Economy." \textit{A.E.R.} (Maio 1967), pp. 534–44. (a)

———. "Inflation and Economic Growth." \textit{J.P.E.} 75 (Dezembro 1967): 796–810. (b)

Snyder, Carl. "On the Statistical Relation of Trade, Credit and Prices." \textit{Rev. Inst. Internat. Statis.} (Outubro 1934), pp. 278–91.

Solow, R. M. "A Contribution to the Theory of Economic Growth." \textit{Q.J.E.} (Fevereiro 1956), pp. 65–94.

———. "On a Family of Lag Distributions." \textit{Econometrica} (Abril 1960), pp. 393–406.

Stein, J. L. "Money and Capacity Growth." \textit{J.P.E.} 74 (Outubro 1966): 451–65.

Tobin, James. "Money Wage Rates and Employment." In \textit{The New Economics}, editado por Seymour Harris. Nova York: Knopf, 1947.

———. "The Interest Elasticity of Transactions Demand for Cash." \textit{Rev. Econ. and Statis.} (Agosto 1956), pp. 241–47.

———. "Liquidity Preference as Behavior toward Risk." \textit{Rev. Econ. Studies} 25 (Fevereiro 1958): 65–86.

———. "The Monetary Interpretation of History." \textit{A.E.R.} (Junho 1965), pp. 464–85. (a)

———. "Money and Economic Growth." \textit{Econometrica} (Outubro 1965), pp. 671–84. (b)

Tobin, James, e Brainard, William C. "Financial Intermediaries and the Effectiveness of Monetary Controls." In \textit{Financial Markets and Economic Activity}. Monografia da Fundação Cowles 21. Nova York: Wiley, 1967. Reimpresso de \textit{A.E.R.} (Maio 1963), pp. 383–400.

———. "Pitfalls in Financial Model Building." \textit{A.E.R.} (Maio 1968), pp. 99–122.

Uzawa, H. "On a Neo-classical Model of Economic Growth." \textit{Econ. Studies Q.} (Setembro 1966), pp. 1–14.

\end{document}